\documentclass[11pt,letterpaper]{article}

\usepackage[margin=1in]{geometry}
\usepackage{fontspec}
\setmainfont{Latin Modern Roman}
\newfontfamily\koreanfont[
  Path=fonts/,
  UprightFont=NotoSansKR-Reconstructed-Regular.ttf,
  BoldFont=NotoSansKR-Reconstructed-Bold.ttf,
  Ligatures=TeX
]{Noto Sans KR}
\newenvironment{koreantranslation}
  {\koreanfont\XeTeXlinebreaklocale "ko"\XeTeXlinebreakskip=0pt plus 1pt}
  {}
\usepackage{microtype}
\usepackage{amsmath,amssymb,mathtools,bm}
\usepackage{booktabs,longtable,array}
\usepackage{tikz}
\usepackage{hyperref}

\hypersetup{
  colorlinks=true,
  linkcolor=black,
  citecolor=black,
  urlcolor=black
}

\setlength{\parindent}{1.5em}
\setlength{\parskip}{0pt}
\allowdisplaybreaks

\newcommand{\epsLJ}{\varepsilon_{\mathrm{LJ}}}
\newcommand{\sigLJ}{\sigma_{\mathrm{LJ}}}
\newcommand{\Zsh}{\mathcal Z}
\newcommand{\Bmn}{\mathcal B}
\newcommand{\dd}{\,\mathrm d}
\newcommand{\ii}{\mathrm i}
\newcommand{\kB}{k_{\mathrm B}}

\title{\vspace{-2.0em}
One-Dimensional Fatigue Modeling of Aluminum:\\
Exact General-$m,n$ Lattice Energy, Inertial Probability Dynamics,\\
and Plastic-Slip Extension}
\author{Donghyeok Kim \and Minsoo Kang \and Junhyeok Park}
\date{}

\begin{document}

\maketitle

\begin{abstract}
This note derives, without fixing the repulsive and attractive exponents, the exact interaction
energy of two parallel infinite atomic rows separated by a normal distance $a$ and shifted by a
tangential registry $s$. A generalized Lennard--Jones $m$--$n$ pair interaction produces a shifted
inhomogeneous Epstein--Hurwitz sum. Mellin transformation and Poisson summation give an exact
Fourier--Bessel representation; no Taylor expansion of the energy landscape is used. The
derivation is carried through to the excess slip energy, internal traction, representative-patch
energy, equilibrium density, generalized Langevin and Kramers dynamics, the conditional
Smoluchowski limit, moment equations, and irreversible dissipation. The force--velocity mobility
law is not used as an atomistic first principle: it appears only after explicit elimination of
rapid inertial motion. An unwrapped well index converts periodic registry changes into a reduced plastic
slip variable. The model can therefore represent ideal thermally activated or instability-driven
single slip and its hysteresis, while quantitative plasticity of face-centered-cubic aluminum
requires a crystallographic half-space energy obtained from a validated EAM or first-principles
model.
\end{abstract}

\section{Definitions, Counting Convention, and Scope}

Let $a>0$ be a normal atomic or interplanar spacing, $s$ a relative tangential registry, and $b$
the lattice repeat distance along a prescribed slip direction. The repulsive and attractive
exponents remain symbolic throughout:
\begin{equation}
  \boxed{m>n>1.}
  \label{eq:exponent_condition}
\end{equation}
The lower bound on $n$ is not a fitted choice: it is the convergence condition for an infinite
one-dimensional lattice sum. The scalar $s$ is an internal microscopic coordinate. Introducing it
does not require a two-dimensional finite-element constitutive law; the surrounding finite-element
mesh may still be one-, two-, or three-dimensional while the present material law uses only a
selected normal component and, if plastic slip is activated, one resolved slip coordinate.

For a finite set of atoms the pair energy is counted once as
\begin{equation}
  E_{\rm pair}=\frac12\sum_i\sum_{j\ne i}v(r_{ij})
              =\sum_{i<j}v(r_{ij}).
  \label{eq:pair_counting}
\end{equation}
The factor one half is essential when both directions from every atom are summed. By contrast, an
interaction between one chosen upper atom and all lower-row atoms contains no factor one half,
because every cross-row pair in that local sum occurs only once.

\section{Generalized Lennard--Jones \texorpdfstring{$m$--$n$}{m--n} Interaction}

It is convenient to preserve the usual meanings of $\epsLJ$ and $\sigLJ$: $\epsLJ$ is the well
depth and $\sigLJ$ is the zero-crossing distance. The normalized general-$m,n$ potential is
\begin{equation}
  \boxed{
  v_{m,n}(r)=C_{m,n}\epsLJ
  \left[
    \left(\frac{\sigLJ}{r}\right)^m
    -\left(\frac{\sigLJ}{r}\right)^n
  \right],
  \qquad
  C_{m,n}=\frac{m}{m-n}
  \left(\frac{m}{n}\right)^{\!n/(m-n)}.}
  \label{eq:generalized_LJ}
\end{equation}
To derive the normalization, set $v'_{m,n}(r_e)=0$. Differentiation gives
\begin{equation}
  -m\sigLJ^m r_e^{-m-1}+n\sigLJ^n r_e^{-n-1}=0,
\end{equation}
hence
\begin{equation}
  \boxed{r_e=\sigLJ\left(\frac{m}{n}\right)^{1/(m-n)}.}
  \label{eq:re_sigma}
\end{equation}
Requiring $v_{m,n}(r_e)=-\epsLJ$ then yields the coefficient $C_{m,n}$ in
Eq.~\eqref{eq:generalized_LJ}. Eliminating $\sigLJ$ by Eq.~\eqref{eq:re_sigma} gives the exactly
equivalent equilibrium-distance form
\begin{equation}
  \boxed{
  v_{m,n}(r)=\epsLJ\left[
    \frac{n}{m-n}\left(\frac{r_e}{r}\right)^m
    -\frac{m}{m-n}\left(\frac{r_e}{r}\right)^n
  \right].}
  \label{eq:LJ_re_form}
\end{equation}
Thus $v(r_e)=-\epsLJ$, $v'(r_e)=0$, and $v''(r_e)>0$ follow without assigning numerical values
to either exponent.

\subsection{Exact normal-chain energy and the two-direction issue}

For $N$ equally spaced atoms in a straight finite chain,
\begin{equation}
  E_N(a)=\sum_{i<j}v_{m,n}((j-i)a)
        =\sum_{k=1}^{N-1}(N-k)v_{m,n}(ka).
  \label{eq:finite_chain_energy}
\end{equation}
The factor $N-k$ is the exact number of pairs separated by $ka$. Dividing by $N$ and taking the
thermodynamic limit gives the energy per atom
\begin{align}
  U_\infty(a)
  &=\lim_{N\to\infty}\frac{E_N(a)}{N}
    =\sum_{k=1}^{\infty}v_{m,n}(ka)\notag\\
  &=\boxed{C_{m,n}\epsLJ\left[
    \left(\frac{\sigLJ}{a}\right)^m\zeta(m)
    -\left(\frac{\sigLJ}{a}\right)^n\zeta(n)
  \right].}
  \label{eq:normal_chain_energy}
\end{align}
An interior reference atom interacts in both directions, so its unhalved local interaction is
$2U_\infty$. Multiplication by one half in Eq.~\eqref{eq:pair_counting} returns the thermodynamic
per-atom value $U_\infty$. Equation~\eqref{eq:normal_chain_energy} therefore already includes both
directions with correct pair counting. A terminal atom in a semi-infinite chain has neighbors in
only one direction and happens to have the same sum, but for a different geometrical reason.

The exact zero-load equilibrium of the infinite chain follows from $U_\infty'(a_0)=0$:
\begin{equation}
  \boxed{
  a_0=\sigLJ
  \left[\frac{m\zeta(m)}{n\zeta(n)}\right]^{1/(m-n)}.}
  \label{eq:chain_equilibrium}
\end{equation}
This is a lattice equilibrium and is generally not identical to the isolated-pair distance $r_e$.

\section{Two-Row Geometry and Direct Slip-Dependent Sum}

Let $\bm d$ be a unit vector along the selected registry direction and $\bm n$ the unit normal,
with $\bm d\mathbin{\cdot}\bm n=0$. The lower row is
\begin{equation}
  \bm R_p^-=pb\bm d,\qquad p\in\mathbb Z,
\end{equation}
and the vector from one upper reference atom to the $p$th lower atom is
\begin{equation}
  \bm r_p(a,s)=a\bm n+(pb+s)\bm d.
\end{equation}
Therefore
\begin{equation}
  \boxed{r_p(a,s)=\sqrt{a^2+(pb+s)^2}.}
  \label{eq:pair_distance}
\end{equation}

\begin{figure}[ht]
  \centering
  \begin{tikzpicture}[x=1.05cm,y=1.0cm,>=stealth]
    \foreach \x in {-3,-2,-1,0,1,2,3}
      \fill (\x,0) circle (2.2pt);
    \foreach \x in {-2.65,-1.65,-0.65,0.35,1.35,2.35}
      \draw[fill=white] (\x,1.75) circle (2.2pt);
    \draw[<->] (-3,-0.35)--(-2,-0.35) node[midway,below] {$b$};
    \draw[dashed] (0,0)--(0,1.75);
    \draw[<->] (0.12,0)--(0.12,1.75) node[midway,right] {$a$};
    \draw[<->] (0,2.12)--(-0.65,2.12) node[midway,above] {$s$};
    \draw[->] (3.2,0)--(3.9,0) node[right] {$\bm d$};
    \draw[->] (3.2,0)--(3.2,0.7) node[above] {$\bm n$};
    \node[left] at (-3.2,0) {lower row};
    \node[left] at (-3.2,1.75) {upper row};
  \end{tikzpicture}
  \caption{Reduced two-row geometry. The normal separation is $a$, the row repeat is $b$, and
  $s$ is the tangential registry.}
  \label{fig:two_rows}
\end{figure}

The cross-row interaction of the upper reference atom is
\begin{align}
  W_{m,n}(a,s)
  &=\sum_{p=-\infty}^{\infty}v_{m,n}(r_p)\notag\\
  &=C_{m,n}\epsLJ\sum_{p=-\infty}^{\infty}
  \left[
    \frac{\sigLJ^m}{[a^2+(pb+s)^2]^{m/2}}
    -\frac{\sigLJ^n}{[a^2+(pb+s)^2]^{n/2}}
  \right].
  \label{eq:W_direct}
\end{align}
For large $|p|$, the two absolute terms are respectively $O(|p|^{-m})$ and
$O(|p|^{-n})$. Hence Eq.~\eqref{eq:exponent_condition} makes both series absolutely convergent;
no cancellation between divergent repulsive and attractive sums is being used.

Set
\begin{equation}
  \delta=\frac{s}{b},\qquad \eta=\frac{a}{b},
\end{equation}
and define
\begin{equation}
  \boxed{
  \Zsh_\nu(\delta,\eta)=
  \sum_{p=-\infty}^{\infty}
  \big[(p+\delta)^2+\eta^2\big]^{-\nu},
  \qquad \eta>0,\quad \operatorname{Re}\nu>\frac12.}
  \label{eq:Z_definition}
\end{equation}
Then
\begin{equation}
  \boxed{
  W_{m,n}(a,s)=C_{m,n}\epsLJ\left[
    \left(\frac{\sigLJ}{b}\right)^m\Zsh_{m/2}(\delta,\eta)
    -\left(\frac{\sigLJ}{b}\right)^n\Zsh_{n/2}(\delta,\eta)
  \right].}
  \label{eq:W_zeta}
\end{equation}
Because reindexing $p\mapsto p-1$ leaves the sum unchanged,
$W_{m,n}(a,s+b)=W_{m,n}(a,s)$ exactly.

\section{Exact Mellin--Poisson--Bessel Derivation}

\subsection{Mellin representation and interchange of sum and integral}

For $x>0$ and $\operatorname{Re}\nu>0$,
\begin{equation}
  x^{-\nu}=\frac{1}{\Gamma(\nu)}
  \int_0^\infty t^{\nu-1}e^{-xt}\dd t.
  \label{eq:Mellin}
\end{equation}
This identity follows directly from the defining Gamma integral
\begin{equation}
  \Gamma(\nu)=\int_0^\infty u^{\nu-1}e^{-u}\dd u.
\end{equation}
Setting $u=xt$ gives
$\Gamma(\nu)=x^\nu\int_0^\infty t^{\nu-1}e^{-xt}\dd t$; division by
$x^\nu\Gamma(\nu)$ produces Eq.~\eqref{eq:Mellin}. Thus the Mellin step introduces no
approximation: it rewrites each inverse power as a continuous superposition of Gaussians.
Substitution into every positive term of Eq.~\eqref{eq:Z_definition} gives
\begin{align}
  \Zsh_\nu(\delta,\eta)
  &=\frac{1}{\Gamma(\nu)}
  \sum_{p\in\mathbb Z}\int_0^\infty
  t^{\nu-1}e^{-t[(p+\delta)^2+\eta^2]}\dd t\notag\\
  &=\frac{1}{\Gamma(\nu)}\int_0^\infty
  t^{\nu-1}e^{-\eta^2t}
  \sum_{p\in\mathbb Z}e^{-t(p+\delta)^2}\dd t.
  \label{eq:Mellin_sum}
\end{align}
For real $\nu>1/2$ the integrand is nonnegative, so Tonelli's theorem justifies the interchange.
For complex $\nu$ in the same half-plane, absolute convergence gives the corresponding Fubini
justification.

\subsection{Shifted Gaussian under Poisson summation}

Use the Fourier-transform convention
\begin{equation}
  \widehat g(\ell)=\int_{-\infty}^{\infty}g(x)e^{-2\pi\ii\ell x}\dd x.
\end{equation}
For $g(x)=e^{-t(x+\delta)^2}$, the change of variable $y=x+\delta$ gives
\begin{align}
  \widehat g(\ell)
  &=e^{2\pi\ii\ell\delta}
    \int_{-\infty}^{\infty}e^{-ty^2}e^{-2\pi\ii\ell y}\dd y\notag\\
  &=\sqrt{\frac{\pi}{t}}
    e^{-\pi^2\ell^2/t}e^{2\pi\ii\ell\delta}.
  \label{eq:Gaussian_transform}
\end{align}
The last equality is the exact Fourier transform of a Gaussian. It may be checked by completing
the square,
\begin{equation}
  -ty^2-2\pi\ii\ell y
  =-t\left(y+\frac{\pi\ii\ell}{t}\right)^2
   -\frac{\pi^2\ell^2}{t},
\end{equation}
and using $\int_{-\infty}^{\infty}e^{-ty^2}\dd y=\sqrt{\pi/t}$. The complex shift is permitted
because the Gaussian integrand is entire and decays on the shifted contour.
Poisson summation $\sum_p g(p)=\sum_\ell\widehat g(\ell)$ therefore gives
\begin{equation}
  \boxed{
  \sum_{p\in\mathbb Z}e^{-t(p+\delta)^2}
  =\sqrt{\frac{\pi}{t}}
  \sum_{\ell\in\mathbb Z}
  e^{-\pi^2\ell^2/t}e^{2\pi\ii\ell\delta}.}
  \label{eq:Poisson}
\end{equation}
Inserting Eq.~\eqref{eq:Poisson} into Eq.~\eqref{eq:Mellin_sum} yields
\begin{equation}
  \Zsh_\nu=\frac{\sqrt\pi}{\Gamma(\nu)}
  \sum_{\ell\in\mathbb Z}e^{2\pi\ii\ell\delta}
  \int_0^\infty t^{\nu-3/2}
  \exp\!\left(-\eta^2t-\frac{\pi^2\ell^2}{t}\right)\dd t.
  \label{eq:after_Poisson}
\end{equation}

\subsection{Zero and nonzero reciprocal modes}

For $\ell=0$, set $u=\eta^2t$:
\begin{align}
  I_0
  &=\int_0^\infty t^{\nu-3/2}e^{-\eta^2t}\dd t\notag\\
  &=\eta^{1-2\nu}\int_0^\infty u^{\nu-3/2}e^{-u}\dd u
   =\Gamma\!\left(\nu-\frac12\right)\eta^{1-2\nu}.
  \label{eq:zero_mode}
\end{align}
\subsection{Definition of the modified Bessel function \texorpdfstring{$K_\mu$}{Kmu}}

The function $K_\mu(\chi)$ appearing below is the modified Bessel function of the second kind. The
argument is denoted by $\chi$ here so that it cannot be confused with the integer slip-well index
$z$ introduced later. It is
one of the two independent solutions of
\begin{equation}
  \chi^2y''+\chi y'-(\chi^2+\mu^2)y=0.
  \label{eq:modified_Bessel_ODE}
\end{equation}
For $\chi>0$ it can be defined directly by the convergent integral
\begin{equation}
  \boxed{
  K_\mu(\chi)=\frac12\int_0^\infty
  u^{\mu-1}\exp\!\left[-\frac{\chi}{2}\left(u+\frac1u\right)\right]\dd u.}
  \label{eq:K_definition}
\end{equation}
This definition immediately shows $K_{-\mu}(\chi)=K_\mu(\chi)$ after the substitution $u\mapsto1/u$.
For positive real $\chi$, $K_\mu(\chi)$ is real and positive for real $\mu$, and it is the solution that
decays at large argument. For completeness, it is related to the modified Bessel function of the
first kind by
\begin{equation}
  K_\mu(\chi)=\frac{\pi}{2}
  \frac{I_{-\mu}(\chi)-I_\mu(\chi)}{\sin(\pi\mu)},
  \qquad
  I_\mu(\chi)=\sum_{j=0}^{\infty}
  \frac{(\chi/2)^{2j+\mu}}{j!\,\Gamma(j+\mu+1)},
  \label{eq:K_from_I}
\end{equation}
where integer orders are obtained by continuous limiting values. The integral definition
Eq.~\eqref{eq:K_definition}, rather than the relation in Eq.~\eqref{eq:K_from_I}, is the form used
in the lattice-sum derivation.

For $\ell\ne0$, the required reciprocal-mode integral has the generic form
\begin{equation}
  \mathcal I(\mu,\beta,\gamma)
  :=\int_0^\infty t^{\mu-1}e^{-\beta t-\gamma/t}\dd t,
  \qquad \operatorname{Re}\beta>0,\quad\operatorname{Re}\gamma>0.
  \label{eq:Bessel_generic_integral}
\end{equation}
Set
\begin{equation}
  t=\sqrt{\frac{\gamma}{\beta}}\,u,
  \qquad
  \dd t=\sqrt{\frac{\gamma}{\beta}}\dd u.
\end{equation}
Then the power and exponent transform exactly as
\begin{align}
  t^{\mu-1}\dd t
  &=\left(\frac{\gamma}{\beta}\right)^{\mu/2}u^{\mu-1}\dd u,\\
  \beta t+\frac{\gamma}{t}
  &=\sqrt{\beta\gamma}\left(u+\frac1u\right).
\end{align}
Consequently,
\begin{align}
  \mathcal I(\mu,\beta,\gamma)
  &=\left(\frac{\gamma}{\beta}\right)^{\mu/2}
  \int_0^\infty u^{\mu-1}
  \exp\!\left[-\sqrt{\beta\gamma}\left(u+\frac1u\right)\right]\dd u\notag\\
  &=\boxed{2\left(\frac{\gamma}{\beta}\right)^{\mu/2}
  K_\mu(2\sqrt{\beta\gamma})}.
  \label{eq:Bessel_identity}
\end{align}
The last line is exactly Eq.~\eqref{eq:K_definition} with $\chi=2\sqrt{\beta\gamma}$; it is not a
table lookup or a local approximation. For the present lattice sum,
\begin{equation}
  \mu=\nu-\frac12,\qquad \beta=\eta^2,\qquad
  \gamma=\pi^2\ell^2.
\end{equation}
Therefore
\begin{equation}
  I_\ell=2\left(\frac{\pi|\ell|}{\eta}\right)^{\nu-1/2}
  K_{\nu-1/2}(2\pi|\ell|\eta).
  \label{eq:nonzero_mode}
\end{equation}
The order $\nu-1/2$ and argument $2\pi|\ell|\eta$ are dimensionless. Since $\eta>0$ and
$\ell\ne0$, the argument lies on the positive real axis, where the integral definition is
unambiguous and exponentially convergent.
Because $I_{-\ell}=I_\ell$ and
$e^{2\pi\ii\ell\delta}+e^{-2\pi\ii\ell\delta}=2\cos(2\pi\ell\delta)$,
Eq.~\eqref{eq:after_Poisson} becomes
\begin{equation}
  \boxed{
  \begin{aligned}
  \Zsh_\nu(\delta,\eta)=\frac{\sqrt\pi}{\Gamma(\nu)}\Bigg[
  &\Gamma\!\left(\nu-\frac12\right)\eta^{1-2\nu}\\
  &+4\sum_{\ell=1}^{\infty}\cos(2\pi\ell\delta)
  \left(\frac{\pi\ell}{\eta}\right)^{\nu-1/2}
  K_{\nu-1/2}(2\pi\ell\eta)
  \Bigg].
  \end{aligned}}
  \label{eq:Poisson_Bessel}
\end{equation}
Equation~\eqref{eq:Poisson_Bessel} is an identity, not a long-wavelength, small-slip, or harmonic
approximation. Its reciprocal terms converge exponentially because
$K_\alpha(x)\sim\sqrt{\pi/(2x)}e^{-x}$ as $x\to\infty$.

\section{Closed Full Energy for Symbolic \texorpdfstring{$m,n$}{m,n}}

For any exponent $q>1$, define
\begin{equation}
  \boxed{
  \begin{aligned}
  \Bmn_q(\delta,\eta):={}&
  \frac{\sqrt\pi}{\Gamma(q/2)}\Bigg[
  \Gamma\!\left(\frac{q-1}{2}\right)\eta^{1-q}\\
  &+4\sum_{\ell=1}^{\infty}\cos(2\pi\ell\delta)
  \left(\frac{\pi\ell}{\eta}\right)^{(q-1)/2}
  K_{(q-1)/2}(2\pi\ell\eta)
  \Bigg].
  \end{aligned}}
  \label{eq:Bq}
\end{equation}
Thus $\Bmn_q=\Zsh_{q/2}$. Substitution of Eq.~\eqref{eq:Bq} into the generalized pair energy
gives the requested full, exact, exponent-symbolic expression:
\begin{equation}
  \boxed{
  \begin{aligned}
  W_{m,n}(a,s)=C_{m,n}\epsLJ\Bigg\{
  &\left(\frac{\sigLJ}{b}\right)^m
  \frac{\sqrt\pi}{\Gamma(m/2)}
  \Bigg[
    \Gamma\!\left(\frac{m-1}{2}\right)\eta^{1-m}\\
  &\quad+4\sum_{\ell=1}^{\infty}\cos(2\pi\ell\delta)
    \left(\frac{\pi\ell}{\eta}\right)^{(m-1)/2}
    K_{(m-1)/2}(2\pi\ell\eta)
  \Bigg]\\
  -&\left(\frac{\sigLJ}{b}\right)^n
  \frac{\sqrt\pi}{\Gamma(n/2)}
  \Bigg[
    \Gamma\!\left(\frac{n-1}{2}\right)\eta^{1-n}\\
  &\quad+4\sum_{\ell=1}^{\infty}\cos(2\pi\ell\delta)
    \left(\frac{\pi\ell}{\eta}\right)^{(n-1)/2}
    K_{(n-1)/2}(2\pi\ell\eta)
  \Bigg]\Bigg\},\\[-0.2em]
  &\hspace{10em}\delta=\frac{s}{b},\qquad\eta=\frac{a}{b}.
  \end{aligned}}
  \label{eq:full_W_mn}
\end{equation}
Every factor multiplying $\epsLJ$ in Eq.~\eqref{eq:full_W_mn} is dimensionless, so
$[W_{m,n}]={\rm energy}$. No exponent has been assigned a numerical value.

\section{Excess Slip Energy, Traction, and Exact Barrier}

Let $s_0$ be a chosen perfect registry, $\delta_0=s_0/b$, and let $A_{\rm at}$ denote the
slip-plane area represented by one reference atom. Define
\begin{equation}
  \Delta W_{m,n}(a,s)=W_{m,n}(a,s)-W_{m,n}(a,s_0),
  \qquad
  \boxed{\gamma_{m,n}(a,s)=\frac{\Delta W_{m,n}(a,s)}{A_{\rm at}}.}
  \label{eq:gamma_definition}
\end{equation}
The registry-independent zero Fourier mode cancels. For $q>1$ define
\begin{equation}
  \boxed{
  \begin{aligned}
  \Delta\Bmn_q(\delta,\delta_0;\eta)
  =\frac{4\sqrt\pi}{\Gamma(q/2)}
  \sum_{\ell=1}^{\infty}
  &[\cos(2\pi\ell\delta)-\cos(2\pi\ell\delta_0)]\\
  &\times\left(\frac{\pi\ell}{\eta}\right)^{(q-1)/2}
  K_{(q-1)/2}(2\pi\ell\eta).
  \end{aligned}}
  \label{eq:DeltaBq}
\end{equation}
The complete generalized slip energy is therefore
\begin{equation}
  \boxed{
  \gamma_{m,n}(a,s)=\frac{C_{m,n}\epsLJ}{A_{\rm at}}
  \left[
    \left(\frac{\sigLJ}{b}\right)^m\Delta\Bmn_m
    -\left(\frac{\sigLJ}{b}\right)^n\Delta\Bmn_n
  \right].}
  \label{eq:gamma_mn}
\end{equation}
The internal resisting traction is
\begin{equation}
  \tau_{\rm int}(a,s)=\frac{\partial\gamma_{m,n}}{\partial s}
  =\frac{C_{m,n}\epsLJ}{A_{\rm at}b}
  \left[
    \left(\frac{\sigLJ}{b}\right)^m\partial_\delta\Bmn_m
    -\left(\frac{\sigLJ}{b}\right)^n\partial_\delta\Bmn_n
  \right],
  \label{eq:traction_full}
\end{equation}
where the exact derivative for either $q=m$ or $q=n$ is
\begin{equation}
  \boxed{
  \begin{aligned}
  \partial_\delta\Bmn_q(\delta,\eta)
  =-\frac{8\pi\sqrt\pi}{\Gamma(q/2)}
  \sum_{\ell=1}^{\infty}\ell\sin(2\pi\ell\delta)
  \left(\frac{\pi\ell}{\eta}\right)^{(q-1)/2}
  K_{(q-1)/2}(2\pi\ell\eta).
  \end{aligned}}
  \label{eq:Bq_derivative}
\end{equation}
Because $[\gamma]={\rm energy}/{\rm area}$ and $[s]={\rm length}$,
$[\tau_{\rm int}]={\rm force}/{\rm area}$. At fixed $a$, the ideal lattice resistance is
\begin{equation}
  \tau_{\rm ideal}(a)=\max_{s\in[s_0,s_0+b)}|\tau_{\rm int}(a,s)|.
  \label{eq:ideal_strength}
\end{equation}

Under a resolved external shear $\tau(t)$, the tilted energy of a patch of area $A_0$ is
\begin{equation}
  \psi(a,s,t)=A_0\gamma_{m,n}(a,s)-\tau(t)A_0(s-s_0).
  \label{eq:tilted_energy}
\end{equation}
A stable registry $s_z$ and a neighboring saddle $s_z^\ddagger$ are obtained from the full
energy, not from a local polynomial:
\begin{align}
  \partial_s\gamma_{m,n}(a,s_z)&=\tau,
  &\partial_s^2\gamma_{m,n}(a,s_z)&>0,\label{eq:stable_well}\\
  \partial_s\gamma_{m,n}(a,s_z^\ddagger)&=\tau,
  &\partial_s^2\gamma_{m,n}(a,s_z^\ddagger)&<0.\label{eq:saddle}
\end{align}
The exact forward or backward activation barrier is
\begin{equation}
  \boxed{
  \Delta G_{z,\pm}^\ddagger(a,\tau)=A_0\left[
  \gamma_{m,n}(a,s_{z,\pm}^\ddagger)-\gamma_{m,n}(a,s_z)
  -\tau(s_{z,\pm}^\ddagger-s_z)
  \right].}
  \label{eq:exact_barrier}
\end{equation}
Loss of metastability occurs when a minimum and saddle coalesce, equivalently
\begin{equation}
  \partial_s\gamma_{m,n}=\tau,
  \qquad
  \partial_s^2\gamma_{m,n}=0.
  \label{eq:loss_stability}
\end{equation}

\section{Full Representative-Patch Energy Under External Loading}

Let a representative patch have area $A_0$ and contain
$N_A=A_0/A_{\rm at}$ equivalent atomic columns. Using the normal chain energy as the reference
and adding only the excess registry energy prevents double counting. The reduced full potential is
\begin{equation}
  \boxed{
  \begin{aligned}
  \Phi(a,s,t)={}&N_A C_{m,n}\epsLJ\left[
    \left(\frac{\sigLJ}{a}\right)^m\zeta(m)
    -\left(\frac{\sigLJ}{a}\right)^n\zeta(n)
  \right]\\
  &+A_0\gamma_{m,n}(a,s)
  -\sigma(t)A_0(a-a_0)
  -\tau(t)A_0(s-s_0).
  \end{aligned}}
  \label{eq:full_patch_energy}
\end{equation}
All four terms have units of energy. The stress-to-atomic-force conversion is not arbitrary:
$F_n=\sigma A_0$ and $F_s=\tau A_0$ follow from the definition of traction on the represented
area. For uniaxial loading and one specified slip system,
\begin{equation}
  \tau(t)=\mathcal M\sigma(t),
  \label{eq:Schmid}
\end{equation}
where $\mathcal M$ is the crystallographic resolved-shear projection. In the strict normal-only
branch one sets $s=s_0$ and $\tau=0$, and Eq.~\eqref{eq:full_patch_energy} reduces to a scalar
normal-spacing energy.

At fixed load and inside a bounded bonded basin $\Omega_b$, statistical mechanics gives
\begin{equation}
  \boxed{
  P_{\rm eq}(a,s\mid\sigma,\tau)=\frac{1}{\mathcal Z_{\rm conf}}
  \exp[-\beta_T\Phi(a,s)],
  \qquad \beta_T=(\kB T)^{-1},}
  \label{eq:equilibrium_density}
\end{equation}
with
\begin{equation}
  \mathcal Z_{\rm conf}=\iint_{\Omega_b}e^{-\beta_T\Phi(a,s)}\dd a\dd s.
\end{equation}
Under dead tensile stress, $-\sigma A_0a$ makes the potential unbounded below as
$a\to\infty$. Consequently there is no globally normalizable stress-controlled equilibrium on
$0<a<\infty$. The physically relevant object before fracture is a metastable bonded-basin density
with an absorbing escape boundary, or a displacement-controlled finite-system ensemble. This
distinction is necessary when the spacing tail is used as a crack-initiation measure.

The normal part of the reduced energy is
\begin{equation}
  \Phi_n(a,t)=N_AU_\infty(a)-\sigma(t)A_0(a-a_0).
\end{equation}
At fixed tensile stress below loss of stability, the outer barrier $a^\ddagger(\sigma)$ is defined
without an arbitrary critical distance by
\begin{equation}
  \boxed{
  N_AU_\infty'(a^\ddagger)=\sigma A_0,
  \qquad U_\infty''(a^\ddagger)<0.}
  \label{eq:normal_barrier}
\end{equation}
The stable spacing obeys the same force balance with positive curvature. The two stationary points
coalesce at the exact lattice inflection spacing
\begin{equation}
  \boxed{
  a_{\rm inf}=\sigLJ\left[
  \frac{m(m+1)\zeta(m)}{n(n+1)\zeta(n)}
  \right]^{1/(m-n)},
  \qquad
  \sigma_{\rm c}=\frac{N_A}{A_0}U_\infty'(a_{\rm inf}).}
  \label{eq:normal_instability}
\end{equation}
Thus a physically derived spacing-tail measure is
\begin{equation}
  \boxed{
  p_{\rm tail}(t;\sigma)=
  \int_{a^\ddagger(\sigma)}^\infty\int P(a,s,t)\dd s\dd a,}
  \label{eq:tail_probability}
\end{equation}
or, with an absorbing boundary at $a^\ddagger$, the accumulated outward probability flux. This
connects crack initiation to the full loss-of-stability landscape rather than to a fitted $a_c$.

\section{Nonequilibrium Dynamics: Inertial Model and Overdamped Reduction}

\subsection{Primary coarse-grained dynamics with inertia and memory}

Force produces acceleration at the atomistic level. A force--velocity relation is therefore not
introduced as a fundamental law. Let $q_a=a$, $q_s=s$, $v_i=\dot q_i$, and
$i\in\{a,s\}$. Projection of all unresolved atomic coordinates onto these two collective
coordinates gives the generalized Langevin form
\begin{equation}
  \boxed{
  m_i^*\ddot q_i(t)=-\partial_{q_i}\Phi
  -\int_0^t\mathcal K_i(t-t')\dot q_i(t')\dd t'+R_i(t),
  \qquad i\in\{a,s\}.}
  \label{eq:GLE}
\end{equation}
Here $m_i^*$ is the effective mass of the collective coordinate, $\mathcal K_i$ is a causal
memory-friction kernel, and $R_i$ is the projected random force. At thermal equilibrium their
fluctuation--dissipation relation is
\begin{equation}
  \boxed{
  \langle R_i(t)\rangle=0,
  \qquad
  \langle R_i(t)R_j(t')\rangle
  =\kB T\,\mathcal K_i(|t-t'|)\delta_{ij}.}
  \label{eq:GLE_FDT}
\end{equation}
Equation~\eqref{eq:GLE} is a reduced description of conservative many-atom dynamics: energy lost
from $a$ or $s$ is transferred to unresolved phonon and atomic modes. A static lattice energy
alone cannot determine $m_i^*$ or $\mathcal K_i$.

If the memory time is short on the resolved time scale, the convolution is replaced by
$\Gamma_i v_i$. The Markovian underdamped Langevin equation is then
\begin{equation}
  \boxed{
  \dot q_i=v_i,
  \qquad
  m_i^*\dot v_i=-\partial_{q_i}\Phi-\Gamma_i v_i
  +\sqrt{2\Gamma_i\kB T}\,\xi_i(t),}
  \label{eq:underdamped_Langevin}
\end{equation}
where $\langle\xi_i(t)\xi_j(t')\rangle=\delta_{ij}\delta(t-t')$. The corresponding phase-space
density $\rho(a,s,v_a,v_s,t)$ obeys the Kramers equation
\begin{equation}
  \boxed{
  \begin{aligned}
  \partial_t\rho={}&-\sum_{i\in\{a,s\}}\partial_{q_i}(v_i\rho)\\
  &+\sum_{i\in\{a,s\}}\frac{1}{m_i^*}\partial_{v_i}
  \left[(\partial_{q_i}\Phi+\Gamma_i v_i)\rho\right]
  +\sum_{i\in\{a,s\}}\frac{\Gamma_i\kB T}{(m_i^*)^2}\partial_{v_i}^2\rho .
  \end{aligned}}
  \label{eq:Kramers_PDE}
\end{equation}
The spacing--registry density requested by the fatigue model is the velocity marginal
\begin{equation}
  \boxed{P(a,s,t)=\iint_{\mathbb R^2}\rho(a,s,v_a,v_s,t)\dd v_a\dd v_s.}
  \label{eq:velocity_marginal}
\end{equation}
In general Eq.~\eqref{eq:velocity_marginal} does not obey a closed equation until a velocity
time-scale reduction is justified.

Define $\dd\Lambda=\dd a\dd s\dd v_a\dd v_s$ and
$\mathcal P_b=\Omega_b\times\mathbb R^2$. The Kramers free-energy functional is
\begin{equation}
  \boxed{
  \mathcal F_K[\rho,t]=\int_{\mathcal P_b}\rho\left[
  \Phi+\frac12m_a^*v_a^2+\frac12m_s^*v_s^2
  +\kB T\ln\!\left(\frac{\rho}{\rho_*}\right)
  \right]\dd\Lambda .}
  \label{eq:Kramers_free_energy}
\end{equation}
The irreversible velocity currents are
\begin{equation}
  \boxed{
  \mathcal J_{v_i}^{\rm irr}
  =-\frac{\Gamma_i}{m_i^*}v_i\rho
  -\frac{\Gamma_i\kB T}{(m_i^*)^2}\partial_{v_i}\rho .}
  \label{eq:velocity_irreversible_current}
\end{equation}
For closed or sufficiently decaying phase-space boundaries, direct integration by parts yields
\begin{equation}
  \boxed{
  \frac{\dd\mathcal F_K}{\dd t}
  =\langle\partial_t\Phi\rangle_\rho-\dot D_K,
  \qquad
  \dot D_K=\sum_{i\in\{a,s\}}\int_{\mathcal P_b}
  \frac{(m_i^*)^2}{\Gamma_i\rho}
  \left(\mathcal J_{v_i}^{\rm irr}\right)^2\dd\Lambda\geq0.}
  \label{eq:Kramers_dissipation}
\end{equation}
This is the inertial Markov model's nonnegative dissipation law. With a non-Markovian memory
kernel, dissipation is history dependent and Eq.~\eqref{eq:Kramers_dissipation} is replaced by the
corresponding bath-plus-system energy balance.

\subsection{Conditions for the Smoluchowski reduction}

Let $\omega_{\rm L}=2\pi f$ be the loading angular frequency and let $\tau_{K_i}$ characterize
the decay time of $\mathcal K_i$. The overdamped reduction requires, at minimum,
\begin{equation}
  \boxed{
  \epsilon_i^{\rm in}=\frac{\omega_{\rm L}m_i^*}{\Gamma_i}\ll1,
  \qquad
  \epsilon_i^{\rm mem}=\omega_{\rm L}\tau_{K_i}\ll1.}
  \label{eq:overdamped_conditions}
\end{equation}
Only under this verified time-scale separation may the inertial term be eliminated and
\begin{equation}
  \boxed{M_i=\Gamma_i^{-1},\qquad D_i=\kB T M_i}
  \label{eq:Einstein_relation}
\end{equation}
be introduced as the long-time mobility and configuration-space diffusivity. The reduced
overdamped Langevin equations are
\begin{align}
  \dot a&=-M_a\partial_a\Phi+\sqrt{2D_a}\,\xi_a(t),\\
  \dot s&=-M_s\partial_s\Phi+\sqrt{2D_s}\,\xi_s(t).
\end{align}
Their joint Smoluchowski equation is
\begin{equation}
  \boxed{
  \frac{\partial P}{\partial t}=-\frac{\partial J_a}{\partial a}
  -\frac{\partial J_s}{\partial s},}
  \label{eq:Smoluchowski}
\end{equation}
where
\begin{align}
  J_a&=-M_a\left(P\partial_a\Phi+\kB T\partial_aP\right),
  \label{eq:Ja}\\
  J_s&=-M_s\left(P\partial_s\Phi+\kB T\partial_sP\right).
  \label{eq:Js}
\end{align}
Here $[M_i]={\rm length}^2/({\rm energy}\,{\rm time})$,
$[P]={\rm length}^{-2}$, and $[J_i]=({\rm length}\,{\rm time})^{-1}$. These equations describe
terminal drift after fast inertial relaxation; they do not assert that atomic velocity is
fundamentally proportional to force.

\subsection{Normalization, moments, boundaries, and crack probability}

For a closed intact domain with no-flux boundaries,
\begin{equation}
  \boxed{\iint_{\Omega_b}P(a,s,t)\dd a\dd s=1.}
  \label{eq:normalization}
\end{equation}
With an absorbing fracture boundary, the same integral is the intact survival probability
\begin{equation}
  \boxed{S(t)=\iint_{\Omega_b}P\dd a\dd s,\qquad
  P_{\rm crack}(t)=1-S(t)
  =\int_0^t\int_{\partial\Omega_{\rm fr}}\bm J\mathbin{\cdot}\bm n_\Omega\dd S\dd t'.}
  \label{eq:survival_flux}
\end{equation}
The last equality assumes a normalized initial intact population and no other probability loss.
The principal unconditioned moments are
\begin{align}
  \bar a(t)&=\iint_{\Omega_b}aP\dd a\dd s,
  \label{eq:mean_a}\\
  \bar s(t)&=\iint_{\Omega_b}sP\dd a\dd s,
  \label{eq:mean_s}\\
  \operatorname{Var}[a](t)&=\iint_{\Omega_b}[a-\bar a(t)]^2P\dd a\dd s,
  \label{eq:var_a}\\
  \bar E(t)&=\iint_{\Omega_b}\Phi(a,s,t)P\dd a\dd s.
  \label{eq:mean_energy}
\end{align}
For an absorbing problem these are moments weighted by survival; conditional moments among intact
states are divided by $S(t)$. For any smooth $g(a,s)$, the exact boundary-aware identity is
\begin{equation}
  \boxed{
  \frac{\dd\langle g\rangle}{\dd t}
  =\iint_{\Omega_b}\bm\nabla g\mathbin{\cdot}\bm J\dd a\dd s
  -\int_{\partial\Omega_b}g\,\bm J\mathbin{\cdot}\bm n_\Omega\dd S.}
  \label{eq:general_moment_boundary}
\end{equation}
Thus, for no-flux boundaries or sufficiently rapid decay at infinity and constant $M_a$,
\begin{align}
  \frac{\dd\bar a}{\dd t}
  &=\iint J_a\dd a\dd s
   =-M_a\left\langle\partial_a\Phi\right\rangle,
  \label{eq:mean_a_evolution}\\
  \frac{\dd\langle a^2\rangle}{\dd t}
  &=-2M_a\left\langle a\partial_a\Phi\right\rangle+2M_a\kB T.
  \label{eq:second_moment}
\end{align}
These equations are exact in the overdamped model but not generally closed. A quadratic moment
closure would discard the full tails and barriers, so the numerical model advances $P$ itself.

\subsection{Overdamped free-energy balance and its boundary correction}

The nonequilibrium configuration-space free-energy functional and its variational derivative are
\begin{align}
  \mathcal F[P,t]&=\iint_{\Omega_b}
  \left[\Phi P+\kB T P\ln\!\left(\frac{P}{P_*}\right)\right]\dd a\dd s,
  \label{eq:free_energy_functional}\\
  \mu_P&=\frac{\delta\mathcal F}{\delta P}
  =\Phi+\kB T\left[\ln\!\left(\frac{P}{P_*}\right)+1\right].
  \label{eq:probability_chemical_potential}
\end{align}
Consequently $J_i=-M_iP\partial_{q_i}\mu_P$. Integration by parts gives the boundary-aware
identity
\begin{equation}
  \boxed{
  \frac{\dd\mathcal F}{\dd t}
  =\left\langle\partial_t\Phi\right\rangle-\dot D_{\rm irr}
  -\int_{\partial\Omega_b}\mu_P\bm J\mathbin{\cdot}\bm n_\Omega\dd S,}
  \label{eq:dissipation_with_boundary}
\end{equation}
where
\begin{equation}
  \boxed{
  \dot D_{\rm irr}=\iint_{\Omega_b}
  \left[\frac{J_a^2}{M_aP}+\frac{J_s^2}{M_sP}\right]\dd a\dd s
  =\iint_{\Omega_b}\sum_{i\in\{a,s\}}M_iP(\partial_{q_i}\mu_P)^2\dd a\dd s
  \geq0.}
  \label{eq:dissipation}
\end{equation}
The familiar reduced balance
\begin{equation}
  \boxed{
  \frac{\dd\mathcal F}{\dd t}
  =\left\langle\partial_t\Phi\right\rangle-\dot D_{\rm irr}}
  \label{eq:dissipation_closed}
\end{equation}
therefore requires no-flux or periodic boundaries. It must not be applied unchanged to an
absorbing fracture boundary. The second form of Eq.~\eqref{eq:dissipation} is preferred
analytically and numerically near small $P$ because it avoids an apparent division by zero.

The accumulated irreversible energy in the overdamped reduced model is
\begin{equation}
  \boxed{D_{\rm irr}(t)=\int_0^t\dot D_{\rm irr}(t')\dd t'.}
  \label{eq:accumulated_dissipation}
\end{equation}
Since
\begin{equation}
  \left\langle\partial_t\Phi\right\rangle
  =-A_0\dot\sigma(\bar a-a_0)-A_0\dot\tau(\bar s-s_0),
  \label{eq:external_power}
\end{equation}
integration by parts over a closed periodic steady-state cycle gives
\begin{equation}
  \boxed{
  \oint\left[\sigma A_0\dd\bar a+\tau A_0\dd\bar s\right]
  =\int_{\rm cycle}\dot D_{\rm irr}\dd t\geq0.}
  \label{eq:hysteresis_loop}
\end{equation}
Finite relaxation creates dynamic loading--unloading lag; interwell transitions create permanent
reduced slip. In the strict normal-only branch, $s=s_0$, $J_s=0$, and only the $a$ or
$(a,v_a)$ equations remain. Permanent crystallographic plasticity is not claimed in that branch.

\section{Can the Model Represent Plastic Deformation?}

\subsection{Reduced ideal slip: yes, if the registry is unwrapped}

The energy $\gamma_{m,n}(a,s)$ is periodic, so storing only $s\bmod b$ loses the number of completed
lattice translations. Write
\begin{equation}
  \boxed{s(t)=z(t)b+\widetilde s(t),
  \qquad z\in\mathbb Z,\quad 0\leq\widetilde s<b.}
  \label{eq:unwrapped_slip}
\end{equation}
Motion of $\widetilde s$ within one well is predominantly recoverable. A change in the integer
$z$ is an irreversible slip event at the reduced-model level. If $h_{\rm slip}$ is the thickness
over which the displacement discontinuity is homogenized, the plastic shear strain is
\begin{equation}
  \boxed{\gamma_p(t)=\frac{b}{h_{\rm slip}}\langle z(t)\rangle,}
  \qquad
  \varepsilon_p(t)=\mathcal M\gamma_p(t)
  \label{eq:plastic_strain}
\end{equation}
for a single slip system under a uniaxial projection $\mathcal M$. This supplies a permanent
strain even after $\widetilde s$ returns near its original position.

\subsection{Exact overdamped fixed-load first-passage rate, without a local energy expansion}

At fixed $a$ and frozen load, let $s_L$ be a reflecting boundary on the left of a well and $s_R$
an absorbing boundary beyond the selected saddle. The mean first-passage time $T_{\rm FP}(s)$ of
the slip Smoluchowski process obeys the backward equation only after
Eq.~\eqref{eq:overdamped_conditions} has been verified:
\begin{equation}
  -1=-M_s\psi'(s)T_{\rm FP}'(s)+M_s\kB T T_{\rm FP}''(s),
  \qquad T_{\rm FP}'(s_L)=0,\quad T_{\rm FP}(s_R)=0.
  \label{eq:backward_MFPT}
\end{equation}
Multiplying by $e^{-\beta_T\psi}$, integrating once from $s_L$ to an intermediate coordinate,
and then applying the absorbing condition gives the exact quadrature
\begin{equation}
  \boxed{
  T_{\rm FP}(s)=\frac{1}{M_s\kB T}
  \int_s^{s_R}e^{\beta_T\psi(y)}
  \left[\int_{s_L}^{y}e^{-\beta_T\psi(x)}\dd x\right]\dd y.}
  \label{eq:exact_MFPT}
\end{equation}
The frozen-load effective escape rate can be defined by
\begin{equation}
  \boxed{k_{z,+}^{\rm FP}(a,\tau)=\left[T_{\rm FP}(s_z)\right]^{-1},}
  \label{eq:exact_FP_rate}
\end{equation}
with the boundaries reversed for the backward rate. Equation~\eqref{eq:exact_MFPT} is exact for
the stated one-dimensional, frozen-load, overdamped boundary-value problem. It evaluates the full
energy $\psi$ and makes no quadratic or Taylor approximation, but it is not the exact
first-passage formula for underdamped Kramers dynamics. A time-local master equation
using this rate additionally assumes separation between intrawell equilibration, interwell escape,
and load variation. If that separation fails, the time-dependent Smoluchowski equation
\eqref{eq:Smoluchowski}, or Eq.~\eqref{eq:Kramers_PDE} when inertia matters, should be solved
directly.

For comparison only, a high-barrier asymptotic reduction of the exact first-passage result is the
overdamped Kramers formula
\begin{equation}
  \boxed{
  k_{z,\pm}(a,t)=\frac{M_s}{2\pi}
  \sqrt{\psi''(a,s_z,t)\left|\psi''(a,s_{z,\pm}^\ddagger,t)\right|}
  \exp[-\beta_T\Delta G_{z,\pm}^\ddagger(a,\tau(t))].}
  \label{eq:Kramers_rate}
\end{equation}
The barrier in this expression is still evaluated from the full energy, but the curvature
prefactor is a local harmonic approximation. It is therefore optional and is not used to define
the primary Kramers PDE or its explicitly conditional Smoluchowski reduction.

With well probabilities $p_z(a,t)$ and either exact first-passage rates under time-scale separation
or explicitly declared asymptotic rates, a discrete--continuous alternative is
\begin{equation}
  \boxed{
  \begin{aligned}
  \partial_t p_z={}&-\partial_aJ_{a,z}
  +k_{z-1,+}p_{z-1}+k_{z+1,-}p_{z+1}
  -(k_{z,+}+k_{z,-})p_z.
  \end{aligned}}
  \label{eq:master_equation}
\end{equation}
It gives
\begin{equation}
  \frac{\dd\langle z\rangle}{\dd t}
  =\sum_z\int[k_{z,+}(a,t)-k_{z,-}(a,t)]p_z(a,t)\dd a,
  \label{eq:mean_z_rate}
\end{equation}
and the plastic contribution to the loop dissipation is $\oint\tau\dd\gamma_p$. Interwell jumps,
finite kinetic relaxation, and an evolving population among $z$ states are the precise mechanisms
that keep unloading from retracing loading.

\subsection{What the reduced model does not yet contain}

The two-row central-force model can represent an ideal periodic lattice resistance, thermally
activated registry jumps, permanent single-slip strain, and dissipative hysteresis. It does not by
itself contain dislocation lines, their elastic fields, multiplication, forest hardening,
backstress, junctions, grain boundaries, or multiple FCC slip systems. Accordingly:
\begin{itemize}
  \item If $s$ is suppressed and only normal stress and $a$ are retained, the model describes
  elastic opening, metastable decohesion, and damage-tail evolution, not crystallographic plasticity.
  \item If one $s$ and its well index $z$ are retained, it is a reduced ideal single-slip plasticity
  model suitable for mechanism studies and mathematical hysteresis.
  \item Quantitative cyclic plasticity of aluminum requires a validated crystalline energy surface
  and additional internal variables for dislocation storage and hardening.
\end{itemize}

\section{EAM Extension of the Probability Framework}

For a metallic crystal, the pair interaction can be replaced by the embedded-atom method
\cite{DawBaskes}. Its total energy is
\begin{equation}
  \boxed{
  E_{\rm EAM}(a,s)=\frac12\sum_i\sum_{j\ne i}\phi(r_{ij}(a,s))
  +\sum_i F\!\left(\rho_i(a,s)\right),
  \qquad
  \rho_i(a,s)=\sum_{j\ne i}f(r_{ij}(a,s)).}
  \label{eq:EAM_total}
\end{equation}
For the proposed density kernel one may use
\begin{equation}
  \boxed{f(r)=f_0\exp\!\left[-\beta_\rho\left(\frac{r}{r_e}-1\right)\right],}
  \label{eq:density_kernel}
\end{equation}
where $\beta_\rho$ is dimensionless and is distinct from the thermodynamic inverse temperature
$\beta_T$. The embedding term is not obtained by inserting $f(r)$ as another pair potential. The
neighbor contributions must first be summed to form $\rho_i$, and the nonlinear function
$F(\rho_i)$ is then evaluated. This order is what makes EAM many-body.

The corresponding generalized stacking-fault energy for a periodic interface cell is
\begin{equation}
  \boxed{
  \gamma_{\rm EAM}(a,s)=\frac{E_{\rm EAM}(a,s)-E_{\rm EAM}(a,s_0)}{A_{\rm cell}}.}
  \label{eq:EAM_GSF}
\end{equation}
Equations~\eqref{eq:full_patch_energy}--\eqref{eq:master_equation} remain structurally unchanged
after replacing $\gamma_{m,n}$ by $\gamma_{\rm EAM}$. The exponential density sum and a general
nonlinear $F$ do not reduce to the same $m,n$ Epstein sum; they should be evaluated from the exact
neighbor geometry rather than forced into an incorrect zeta form. A validated aluminum EAM can
reproduce metallic cohesive energetics, stacking-fault barriers, and atomistic dislocation motion
far more realistically than pair LJ, but a single collective $s$ still represents only idealized
single slip.

\section{Crystalline Half-Space Generalization}

Let $\mathcal U$ and $\mathcal L$ be upper and lower crystalline half-spaces separated by a chosen
plane. For pair interactions the exact excess energy per periodic interface cell is
\begin{equation}
  \boxed{
  \begin{aligned}
  \gamma_{\rm GSF}(a,s)=\frac{1}{A_{\rm cell}}
  \sum_{i\in\mathcal U_{\rm cell}}\sum_{j\in\mathcal L}
  \Big[&v_{m,n}(|\bm R_{ij}+a\bm n+s\bm d|)\\
       &-v_{m,n}(|\bm R_{ij}+a\bm n+s_0\bm d|)\Big].
  \end{aligned}}
  \label{eq:full_GSF}
\end{equation}
There is no factor one half because every cross-plane pair is counted once. In a true FCC model,
two in-plane lattice indices and multiple normal layers replace the single row index. The same
Mellin--Poisson strategy applies to inverse-power pair terms, but the reciprocal lattice is
two-dimensional. For aluminum, a validated FCC generalized stacking-fault surface from EAM or
first principles should ultimately replace the two-row energy \cite{Lu2000}.

\section{Complete Symbol and Computational Index}

This index is normative for the notation used in this document. Local dummy variables are marked as
such; repeated glyphs with different subscripts are not interchangeable.

\small
\begin{longtable}{>{\raggedright\arraybackslash}p{0.16\textwidth}
                  >{\raggedright\arraybackslash}p{0.34\textwidth}
                  >{\raggedright\arraybackslash}p{0.40\textwidth}}
\toprule
Symbol & Meaning and unit & Definition or calculation method\\
\midrule
\endfirsthead
\toprule
Symbol & Meaning and unit & Definition or calculation method\\
\midrule
\endhead
\bottomrule
\endfoot
$a,a_0$ & instantaneous and zero-load equilibrium normal spacing; length & $a$ is the primary
normal coordinate. Compute $a_0$ from Eq.~\eqref{eq:chain_equilibrium}, or from the minimum of a
validated replacement energy.\\
$a^\ddagger,a_{\rm inf}$ & normal barrier and lattice inflection spacing; length & Bracket all roots
of $N_AU_\infty'(a)=\sigma A_0$, classify by $U_\infty''$, and use the negative-curvature root as
$a^\ddagger$. Equation~\eqref{eq:normal_instability} gives $a_{\rm inf}$.\\
$s,s_0,\widetilde s,b$ & unwrapped registry, reference registry, within-cell registry, and slip
repeat; length & $s=zb+\widetilde s$; obtain $b$ and $s_0$ from the chosen lattice geometry and use
$0\leq\widetilde s<b$.\\
$z,\mathbb Z,\ddagger$ & integer well label, set of integers, and transition-state marker;
dimensionless & Increment or decrement $z$ when unwrapped $s$ crosses a periodic cell boundary.
The dagger is not multiplication; it marks a saddle or activation barrier.\\
$r,r_e,r_p,r_{ij}$ & pair distances; length & $r_e$ is Eq.~\eqref{eq:re_sigma}; $r_p$ is
Eq.~\eqref{eq:pair_distance}; $r_{ij}=|\bm R_i-\bm R_j|$ in the selected atomic geometry.\\
$\bm d,\bm n,\bm e$ & unit slip direction, plane normal, and uniaxial loading direction;
dimensionless & Normalize the crystallographic vectors and verify
$\bm d\mathbin{\cdot}\bm n=0$.\\
$\bm R_p^-,\bm r_p,\bm R_{ij}$ & lower-row site, cross-row separation vector, and reference lattice
vector; length & Construct directly from the lattice basis and integer cell indices.\\
$m,n,q$ & repulsive exponent, attractive exponent, and generic exponent; dimensionless & Require
$m>n>1$ for the present one-dimensional sums; they are neither masses nor Schmid factors.\\
$C_{m,n}$ & generalized LJ normalization; dimensionless & Evaluate
Eq.~\eqref{eq:generalized_LJ}.\\
$\epsLJ,\sigLJ$ & pair well depth and zero-crossing distance; energy and length & Obtain from the
declared pair potential or independent cohesive calibration; do not fit them to the fatigue loop.\\
$\begin{gathered}v_{m,n},E_{\rm pair}\\E_N,U_\infty\end{gathered}$ & pair, finite-system, finite-chain, and infinite-chain-per-atom
energies; energy & Use Eqs.~\eqref{eq:generalized_LJ}, \eqref{eq:pair_counting},
\eqref{eq:finite_chain_energy}, and \eqref{eq:normal_chain_energy}.\\
$W_{m,n},\Delta W_{m,n}$ & cross-row energy and registry excess; energy & Evaluate the direct sum or
Eq.~\eqref{eq:full_W_mn}; subtract the $s_0$ value for $\Delta W$.\\
$\gamma_{m,n},\gamma_{\rm EAM},\gamma_{\rm GSF}$ & slip or generalized stacking-fault energy
density; energy/area & Divide exact interface-cell excess energy by its represented area; see
Eqs.~\eqref{eq:gamma_definition}, \eqref{eq:EAM_GSF}, and \eqref{eq:full_GSF}.\\
$\psi,\Phi,\Phi_n$ & tilted slip energy, full patch potential, and normal-only potential; energy &
Use Eqs.~\eqref{eq:tilted_energy} and \eqref{eq:full_patch_energy}; set $s=s_0,\tau=0$ for
$\Phi_n$.\\
$\Delta G_{z,\pm}^\ddagger$ & forward/backward activation barrier; energy & Locate the full-energy
minimum and adjacent saddle by bracketed root finding, classify by the exact second derivative,
then use Eq.~\eqref{eq:exact_barrier}.\\
$A_{\rm at},A_0,A_{\rm cell}$ & area per atomic column, correlated patch area, and interface-cell
area; area & Compute atomic areas from the declared lattice cell. Choose $A_0=N_AA_{\rm at}$ from
a physical correlation area, not an FEM element area.\\
$N,N_A$ & finite atom count and equivalent column count; integers & $N_A=A_0/A_{\rm at}$ must agree
with the energy counting convention.\\
$h_{\rm slip}$ & thickness homogenizing a slip discontinuity; length & Specify from the selected
slip band or coarse-graining thickness independently of mesh size.\\
$\sigma,\sigma_{\rm c},\tau,\tau_{\rm int},\tau_{\rm ideal}$ & applied normal stress, normal
instability stress, resolved shear, internal traction, and ideal shear strength; stress & Use
Eqs.~\eqref{eq:normal_instability}, \eqref{eq:Schmid}, \eqref{eq:traction_full}, and
\eqref{eq:ideal_strength}.\\
$F_n,F_s,\mathcal M$ & normal force, slip force, and resolved-shear projection; force, force,
dimensionless & $F_n=\sigma A_0$, $F_s=\tau A_0$,
$\mathcal M=(\bm d\mathbin{\cdot}\bm e)(\bm n\mathbin{\cdot}\bm e)$, and
$\tau=\mathcal M\sigma$.\\
$\delta,\delta_0,\eta$ & normalized registry, reference registry, and separation; dimensionless &
$\delta=s/b$, $\delta_0=s_0/b$, $\eta=a/b$.\\
$p,k,\ell,i,j$ & direct-row, pair-separation, reciprocal-mode, and atom/component indices; integers
& Their ranges are stated at each sum. $i,j$ in $\delta_{ij}$ label components, not atoms.\\
$\nu,\mu,\chi$ & Mellin exponent, Bessel order, and Bessel argument; dimensionless & In the lattice
sum, $\nu=q/2$, $\mu=\nu-1/2$, and $\chi=2\pi|\ell|\eta$. The Bessel argument is $\chi$, not the
well index $z$.\\
$t,u,x,y$ in transform integrals & local dummy integration/Fourier variables; $y$ is also the
local dependent variable in the displayed Bessel ODE & Their dimensions follow the displayed
substitution and they have no material-state meaning.\\
$\begin{gathered}g(x),\widehat g(\ell),I_0\\I_\ell,\mathcal I\end{gathered}$ & local Poisson-transform function, its Fourier
transform, zero/nonzero Fourier integrals, and a generic Bessel integral & These are derivation-only
objects defined immediately before the equations in which they occur; evaluate them by the exact
Gamma/Bessel identities shown there.\\
$\Gamma(\cdot),\zeta(\cdot),I_\mu,K_\mu$ & Gamma, Riemann-zeta, modified Bessel first-kind, and
modified Bessel second-kind functions & Use standard special-function libraries; verify $K_\mu$
against Eq.~\eqref{eq:K_definition} when testing.\\
$\Zsh_\nu,\Bmn_q,\Delta\Bmn_q$ & shifted Epstein--Hurwitz sum, closed exponent sum, and registry
difference; dimensionless & Use Eqs.~\eqref{eq:Z_definition}, \eqref{eq:Bq}, and
\eqref{eq:DeltaBq}. Increase the reciprocal cutoff to tolerance and confirm test points against
the direct sum.\\
$\beta,\gamma$ in Eq.~\eqref{eq:Bessel_generic_integral} & positive local integration parameters &
For that derivation only, $\beta=\eta^2$ and $\gamma=\pi^2\ell^2$; they are not inverse temperature
or stacking-fault energy.\\
$\begin{gathered}\pi,e,\ii,\operatorname{Re}\\|\cdot|,\exp,\ln,\sum\\
\int,\lim,\max\end{gathered}$ & circle constant, natural-logarithm base, imaginary
unit, real part, norm, and standard operators & $\exp x=e^x$ and $\ln$ is its inverse. Sums and
integrals use the displayed limits; $\lim$, $\max$, and absolute values have their standard exact
meanings.\\
$!,\bmod,\sim,\propto,\mathbb R$ & factorial, remainder operation, asymptotic/equivalence marker,
proportionality marker, and real-number set & $j!=\Gamma(j+1)$ for nonnegative integer $j$;
$s\bmod b$ retains only the within-cell registry, whereas unwrapped $s$ retains slip history.\\
\end{longtable}

\begin{longtable}{>{\raggedright\arraybackslash}p{0.16\textwidth}
                  >{\raggedright\arraybackslash}p{0.34\textwidth}
                  >{\raggedright\arraybackslash}p{0.40\textwidth}}
\toprule
Dynamic symbol & Meaning and unit & Definition or calculation method\\
\midrule
\endfirsthead
\toprule
Dynamic symbol & Meaning and unit & Definition or calculation method\\
\midrule
\endhead
\bottomrule
\endfoot
$t,f,\omega_{\rm L}$ & time, cyclic frequency, and loading angular frequency; time, 1/time, 1/time &
$\omega_{\rm L}=2\pi f$; for nonsinusoidal loading test every relevant harmonic.\\
$T,\kB,\beta_T$ & absolute temperature, Boltzmann constant, and inverse thermal energy; K,
energy/K, 1/energy & $\beta_T=(\kB T)^{-1}$.\\
$q_i,v_i,m_i^*$ & collective coordinate, velocity, and effective mass; length, length/time, mass &
$q_a=a,q_s=s$. For two rigid patches of masses $M_U,M_L$, use
$m_i^*=M_UM_L/(M_U+M_L)$; use the moving mass when the other patch is fixed.\\
$\mathcal K_i,R_i,\tau_{K_i}$ & memory-friction kernel, projected random force, and kernel decay
time & Obtain $\mathcal K_i(t)=\beta_T\langle R_i(0)R_i(t)\rangle$ from a constrained equilibrium
MD projection and estimate $\tau_{K_i}$ from its decay.\\
$\Gamma_i,M_i,D_i$ & Markov friction, long-time mobility, and configuration diffusivity;
mass/time, length$^2$/(energy time), length$^2$/time & If a Markov plateau exists,
$\Gamma_i=\int_0^\infty\mathcal K_i(t)\dd t$, $M_i=1/\Gamma_i$, and $D_i=\kB TM_i$.\\
$\xi_i,\delta_{ij},\delta(t-t')$ & unit white noise, Kronecker delta, and Dirac delta & Normalize by
$\langle\xi_i(t)\xi_j(t')\rangle=\delta_{ij}\delta(t-t')$. Do not confuse $\delta_{ij}$ with
$\delta=s/b$.\\
$\rho,\rho_*,\mathcal P_b,\dd\Lambda$ & phase density, reference phase density, bonded phase domain,
and phase-volume element & Normalize $\rho$ on $\mathcal P_b=\Omega_b\times\mathbb R^2$;
$\rho_*$ only nondimensionalizes the logarithm.\\
$\mathcal J_{v_i}^{\rm irr}$ & irreversible velocity-space probability current & Compute
Eq.~\eqref{eq:velocity_irreversible_current}; it produces the Markov Kramers dissipation.\\
$\mathcal F_K,\dot D_K$ & inertial nonequilibrium free energy and irreversible loss rate; energy,
energy/time & Use Eqs.~\eqref{eq:Kramers_free_energy} and \eqref{eq:Kramers_dissipation}.\\
$\epsilon_i^{\rm in},\epsilon_i^{\rm mem}$ & inertial and memory reduction tests; dimensionless &
Compute Eq.~\eqref{eq:overdamped_conditions}; use Smoluchowski only when both are small for every
relevant loading harmonic.\\
$P,P_{\rm eq},P_*,\mathcal Z_{\rm conf}$ & configuration density, equilibrium density, reference
density, and configurational partition function & $P$ is the velocity marginal of $\rho$; normalize
using Eq.~\eqref{eq:normalization}; use $P_{\rm eq}$ only in a bounded bonded basin.\\
$\begin{gathered}\Omega_b,\partial\Omega_b\\\partial\Omega_{\rm fr},\bm n_\Omega,\dd S\end{gathered}$ & bonded domain, its
boundary, absorbing fracture portion, outward normal, and boundary measure & Place the normal
fracture boundary at the full-energy saddle $a^\ddagger$, not at a fitted distance.\\
$J_a,J_s,\bm J$ & configuration probability currents; 1/(length time) & Evaluate
Eqs.~\eqref{eq:Ja}--\eqref{eq:Js}; use conservative finite-volume face fluxes.\\
$\mathcal F,\mu_P,D_{\rm irr},\dot D_{\rm irr}$ & overdamped free energy, probability chemical
potential, accumulated dissipation, and rate; energy, energy, energy, energy/time & Use
Eqs.~\eqref{eq:free_energy_functional}, \eqref{eq:probability_chemical_potential},
\eqref{eq:dissipation}, and \eqref{eq:accumulated_dissipation}.\\
$S,P_{\rm crack},p_{\rm tail}$ & intact survival, cumulative crack probability, and instantaneous
tail probability; dimensionless & Use Eqs.~\eqref{eq:survival_flux} and
\eqref{eq:tail_probability}; never identify dissipation directly with crack probability.\\
$\bar a,\bar s,\operatorname{Var}[a],\bar E$ & unconditioned means, variance, and mean potential
energy & Use Eqs.~\eqref{eq:mean_a}--\eqref{eq:mean_energy}; divide by $S$ for intact-conditional
moments.\\
$\begin{gathered}g,\bm\nabla,\langle\cdot\rangle\\\partial_t,\dot{(\cdot)},(\cdot)'\end{gathered}$ & generic observable, gradient,
ensemble average, partial time derivative, total time derivative, and one-variable derivative &
Use Eq.~\eqref{eq:general_moment_boundary}; keep its boundary term for absorbing fracture.\\
\end{longtable}

\begin{longtable}{>{\raggedright\arraybackslash}p{0.16\textwidth}
                  >{\raggedright\arraybackslash}p{0.34\textwidth}
                  >{\raggedright\arraybackslash}p{0.40\textwidth}}
\toprule
Slip/EAM symbol & Meaning and unit & Definition or calculation method\\
\midrule
\endfirsthead
\toprule
Slip/EAM symbol & Meaning and unit & Definition or calculation method\\
\midrule
\endhead
\bottomrule
\endfoot
$s_z,s_{z,\pm}^\ddagger$ & stable registry in well $z$ and adjacent forward/backward saddle; length
& Solve $\partial_s\gamma=\tau$ on each periodic interval and classify with
$\partial_s^2\gamma$.\\
$\gamma_p,\varepsilon_p$ & reduced plastic shear and projected axial plastic strain; dimensionless
& Use Eq.~\eqref{eq:plastic_strain}; this is ideal single-slip bookkeeping, not hardening.\\
$s_L,s_R,x,y,T_{\rm FP}$ & reflecting boundary, absorbing boundary, quadrature variables, and mean
first-passage time; length, length, length, length, time & Select boundaries around the exact well
and saddle and evaluate Eq.~\eqref{eq:exact_MFPT} by adaptive quadrature.\\
$k_{z,\pm}^{\rm FP},k_{z,\pm}$ & first-passage and optional asymptotic transition rates; 1/time &
Invert the exact overdamped MFPT under its assumptions; use Eq.~\eqref{eq:Kramers_rate} only for a
verified high barrier.\\
$p_z,J_{a,z}$ & probability density and normal current in well $z$ & Advance
Eq.~\eqref{eq:master_equation} only when intrawell, escape, and load time scales separate.\\
$\begin{gathered}E_{\rm EAM},\phi,F,\rho_i\\f,f_0,\beta_\rho\end{gathered}$ & EAM energy, pair term, embedding function, host
density, density kernel, amplitude, and decay coefficient & Use a complete validated EAM
parameterization. Sum $f(r_{ij})$ before applying nonlinear $F$. $\beta_\rho$ is not $\beta_T$.\\
$\mathcal U,\mathcal L,\mathcal U_{\rm cell}$ & upper half-crystal, lower half-crystal, and upper
interface cell & Generate from the selected crystal orientation and converge neighbor extent or the
validated EAM cutoff.\\
$M_U,M_L$ & masses of upper and lower representative moving patches; mass & Sum atomic masses in
the correlated patch used for the relative collective coordinate.\\
\end{longtable}
\normalsize

\subsection{Reproducible calculation sequence}

\begin{enumerate}
  \item Declare geometry and counting: $m,n,\epsLJ,\sigLJ,b,A_{\rm at},A_0,s_0$ and whether the
  normal-only or optional slip branch is active.
  \item Evaluate $U_\infty$ and $W_{m,n}$. Increase the Bessel-mode cutoff until the energy and all
  required derivatives are stable at the requested tolerance; verify representative points against
  the direct real-space sum.
  \item Locate minima, saddles, and instability points by bracketed root finding on the full
  analytic energy and classify them with full derivatives. No polynomial energy fit is required.
  \item Obtain $m_i^*$ and $\mathcal K_i$ from atomistic geometry and constrained equilibrium MD.
  Test Eq.~\eqref{eq:overdamped_conditions}; do not introduce $M_i$ unless the reduction is
  supported.
  \item If inertia matters, advance Eq.~\eqref{eq:Kramers_PDE} with a conservative,
  positivity-preserving finite-volume method in $(q_i,v_i)$. If the reduction is justified, advance
  Eq.~\eqref{eq:Smoluchowski} with conservative face fluxes. Check positivity, normalization, and
  time-step and domain convergence in either case.
  \item For crack initiation, place the absorbing normal boundary at the full-energy saddle,
  integrate its outward probability flux, and report $p_{\rm tail}$, $S$, and $P_{\rm crack}$
  separately. Use Eq.~\eqref{eq:dissipation_with_boundary}, not the closed-boundary balance.
  \item Repeat with finer state grids, smaller time steps, larger lattice-sum cutoffs, and larger
  velocity domains. Report convergence before interpreting a fatigue hysteresis loop.
\end{enumerate}


\section{Analytical Consistency Checks}

The symbolic formulas admit checks that do not require inserting any parameter values:
\begin{align}
  \Bmn_q(\delta+1,\eta)&=\Bmn_q(\delta,\eta),
  &\Bmn_q(-\delta,\eta)&=\Bmn_q(\delta,\eta),\\
  \partial_\delta\Bmn_q(0,\eta)&=0,
  &\Delta\Bmn_q(\delta_0,\delta_0;\eta)&=0.
\end{align}
The direct and reciprocal forms must agree for every $q>1$, $\eta>0$, and real $\delta$.
At large $\eta$, the registry-dependent Fourier--Bessel terms vanish exponentially. The normal
Riemann-zeta chain in Eq.~\eqref{eq:normal_chain_energy} is a separate collinear geometry; it must
not be obtained by taking $\eta\to0$ and $\delta\to0$ in the two-row formula, because that limit
creates a coincident cross-row pair and diverges.

\section{Conclusion}

The exact general-$m,n$ result is Eq.~\eqref{eq:full_W_mn}, and the corresponding representative
energy under normal and resolved shear loading is Eq.~\eqref{eq:full_patch_energy}. The primary
reduced kinetics retain acceleration and friction through Eq.~\eqref{eq:Kramers_PDE}; the
Smoluchowski equation is used only after Eq.~\eqref{eq:overdamped_conditions} is verified.
Finite-rate relaxation can then generate dynamic loading--unloading lag, while interwell
transitions of the unwrapped registry generate permanent reduced plastic slip. The nonnegative
dissipation is Eq.~\eqref{eq:Kramers_dissipation} in the inertial Markov model and
Eq.~\eqref{eq:dissipation} in its overdamped reduction. Absorbing fracture boundaries require the
extra boundary term in Eq.~\eqref{eq:dissipation_with_boundary}. Thus the framework can be used as
a mathematically consistent ideal-slip mechanism model without fitting a sinusoidal hysteresis
operator. It should not yet be presented as quantitative aluminum cyclic plasticity until the
two-row LJ landscape is replaced or validated against a full FCC EAM or first-principles
stacking-fault surface and the necessary dislocation hardening variables are introduced.

\begin{thebibliography}{9}

\bibitem{DawBaskes}
M. S. Daw and M. I. Baskes,
``Embedded-atom method: Derivation and application to impurities, surfaces, and other defects in
metals,'' \emph{Physical Review B}, vol. 29, pp. 6443--6453, 1984.
\href{https://doi.org/10.1103/PhysRevB.29.6443}{doi:10.1103/PhysRevB.29.6443}.

\bibitem{Lu2000}
G. Lu, N. Kioussis, V. V. Bulatov, and E. Kaxiras,
``Generalized-stacking-fault energy surface and dislocation properties of aluminum,''
\emph{Physical Review B}, vol. 62, pp. 3099--3108, 2000.
\href{https://doi.org/10.1103/PhysRevB.62.3099}{doi:10.1103/PhysRevB.62.3099}.

\bibitem{DLMFBessel}
NIST Digital Library of Mathematical Functions,
``Integral Representations of Modified Bessel Functions,'' Sec. 10.32.
\url{https://dlmf.nist.gov/10.32}.

\bibitem{DLMFZeta}
NIST Digital Library of Mathematical Functions,
``Riemann Zeta Function: Definition and Expansions,'' Sec. 25.2.
\url{https://dlmf.nist.gov/25.2}.

\end{thebibliography}

\clearpage
\begin{koreantranslation}

\begin{center}
  {\LARGE\bfseries 한국어 번역본}\par
  \vspace{0.5em}
  {\Large 일반 $m,n$ 슬립 의존 격자 에너지와 소성 슬립 확장}\par
  \vspace{1em}
  {\normalsize Donghyeok Kim \quad Minsoo Kang \quad Junhyeok Park}
\end{center}

\vspace{1em}

\noindent\textbf{번역본 안내.}
이 번역부는 앞의 영문 본문을 한국어로 읽을 수 있도록 같은 순서로 다시 설명한다.
수식은 언어와 무관하므로 핵심식은 번역부에도 다시 수록하고, 긴 중간 대수식은 영문
본문의 식 번호를 함께 인용한다. 따라서 한국어 부분만 먼저 읽어도 전체 모델의 구조를
파악할 수 있고, 세부 계산을 검산할 때에는 해당 영문 식으로 바로 돌아갈 수 있다.
여기서 $m,n$에는 특정 숫자를 대입하지 않으며 끝까지 $m>n>1$의 일반형을 유지한다.

\section*{\koreanfont 전체 기호 및 계산 인덱스}

아래 표는 본문의 기호를 빠짐없이 찾기 위한 한국어 인덱스다. 같은 모양이라도
아래첨자가 다르면 다른 양이며, 적분 안에서만 쓰는 dummy variable은 재료변수가 아니다.

\small
\begin{longtable}{>{\raggedright\arraybackslash}p{0.17\textwidth}
                  >{\raggedright\arraybackslash}p{0.34\textwidth}
                  >{\raggedright\arraybackslash}p{0.39\textwidth}}
\toprule
기호 & 의미와 단위 & 정의와 계산방법\\
\midrule
\endfirsthead
\toprule
기호 & 의미와 단위 & 정의와 계산방법\\
\midrule
\endhead
\bottomrule
\endfoot
$a,a_0$ & 현재 및 무하중 평형 법선거리; 길이 & $a$는 핵심 상태변수이고 $a_0$는
식~\eqref{eq:chain_equilibrium} 또는 검증된 에너지의 최소점에서 계산한다.\\
$a^\ddagger,a_{\rm inf}$ & 법선 장벽과 변곡거리; 길이 & 전체 힘평형식의 모든 근을
구간포착한 뒤 곡률 부호로 분류한다. 음의 곡률 근이 $a^\ddagger$이다.\\
$s,s_0,\widetilde s,b$ & 누적 슬립좌표, 기준 registry, 셀 내부좌표, 격자주기; 길이 &
$s=zb+\widetilde s$, $0\leq\widetilde s<b$이며 $b,s_0$는 결정기하에서 정한다.\\
$z,\mathbb Z,\ddagger,\pm$ & 우물번호, 정수집합, 전이상태 표시, 전진/후진; 무차원 &
$z$는 공간좌표가 아니라 장벽통과 누적번호이다. $\ddagger$는 곱셈표시가 아니다.\\
$r,r_e,r_p,r_{ij}$ & 원자쌍 거리; 길이 & $r_e$는 식~\eqref{eq:re_sigma}, $r_p$는
식~\eqref{eq:pair_distance}, $r_{ij}$는 실제 이웃좌표에서 계산한다.\\
$\bm d,\bm n,\bm e$ & 슬립방향, 면법선, 단축하중방향 단위벡터 & 모두 정규화하고
$\bm d\cdot\bm n=0$을 확인한다.\\
$\bm R_p^-,\bm r_p,\bm R_{ij}$ & 격자점 및 원자쌍 기준벡터; 길이 & 선언한 lattice
basis와 정수 cell index로 생성한다.\\
$m,n,q$ & 반발지수, 인력지수, 일반지수; 무차원 & 현재 1D 합의 수렴조건은 $m>n>1$이다.
질량이나 Schmid factor와 무관하다.\\
$C_{m,n},\epsLJ,\sigLJ$ & LJ 정규화계수, 우물깊이, 영점거리; 무차원, 에너지, 길이 &
$C_{m,n}$은 식~\eqref{eq:generalized_LJ}; LJ 계수는 독립적인 응집에너지 자료에서
정하고 피로루프에 맞추지 않는다.\\
$\begin{gathered}v_{m,n},E_{\rm pair}\\E_N,U_\infty\end{gathered}$ & 원자쌍, 유한계, 사슬, 무한사슬 에너지; 에너지 &
각각 식~\eqref{eq:generalized_LJ}, \eqref{eq:pair_counting},
\eqref{eq:finite_chain_energy}, \eqref{eq:normal_chain_energy}로 계산한다.\\
$W_{m,n},\Delta W_{m,n}$ & 두 원자열 전체에너지와 registry 초과에너지; 에너지 &
직접합 또는 식~\eqref{eq:full_W_mn}을 쓰고 기준 $s_0$ 값을 뺀다.\\
$\gamma_{m,n},\gamma_{\rm EAM},\gamma_{\rm GSF}$ & 슬립/GSF 면적에너지; 에너지/면적 &
정확한 interface-cell 초과에너지를 대표면적으로 나눈다.\\
$\psi,\Phi,\Phi_n$ & 기울어진 슬립에너지, 대표영역 전체에너지, normal-only 에너지 &
식~\eqref{eq:tilted_energy}, \eqref{eq:full_patch_energy}; normal-only에서는
$s=s_0,\tau=0$.\\
$\Delta G_{z,\pm}^\ddagger$ & 전진/후진 활성화장벽; 에너지 & 최소점과 인접 안장점을
전체에너지에서 구하고 식~\eqref{eq:exact_barrier}에 대입한다.\\
$A_{\rm at},A_0,A_{\rm cell}$ & 원자기둥 면적, 상관 대표면적, interface-cell 면적 & 결정
단위셀에서 계산하며 $A_0=N_AA_{\rm at}$이다. $A_0$를 FEM 요소면적으로 두지 않는다.\\
$N,N_A,h_{\rm slip}$ & 원자수, 등가 원자기둥수, 슬립 균질화두께 & $N_A=A_0/A_{\rm at}$;
$h_{\rm slip}$은 물리적 슬립밴드/축약두께에서 정한다.\\
$\sigma,\sigma_{\rm c},\tau,\tau_{\rm int},\tau_{\rm ideal}$ & 법선응력, 법선임계응력,
resolved shear, 내부전단저항, 이상전단강도; 응력 & 본문 식
\eqref{eq:normal_instability}, \eqref{eq:Schmid}, \eqref{eq:traction_full},
\eqref{eq:ideal_strength}로 계산한다.\\
$F_n,F_s,\mathcal M$ & 법선 및 접선 일반화힘과 Schmid 투영계수 & $F_n=\sigma A_0$,
$F_s=\tau A_0$, $\mathcal M=(\bm d\cdot\bm e)(\bm n\cdot\bm e)$이다.\\
$\delta,\delta_0,\eta$ & 무차원 registry, 기준 registry, 법선거리 & $\delta=s/b$,
$\delta_0=s_0/b$, $\eta=a/b$.\\
$p,k,\ell,i,j$ & 직접격자, 분리거리, 역격자, 원자/성분 index; 정수 & 각 합에 적힌
범위를 따른다. $\delta_{ij}$의 $i,j$는 성분번호이다.\\
$\nu,\mu,\chi$ & Mellin 지수, Bessel 차수, Bessel 인자 & $\nu=q/2$,
$\mu=\nu-1/2$, $\chi=2\pi|\ell|\eta$. $\chi$를 우물번호 $z$와 분리했다.\\
$t,u,x,y$ & 문맥별 시간 또는 dummy 적분변수; Bessel ODE의 $y$는 국소 종속변수 &
동역학의 $t$ 외에는 재료상태가 아니다.\\
$g(x),\widehat g(\ell),I_0,I_\ell,\mathcal I$ & Poisson 변환함수와 Fourier 변환, 영/비영
모드 적분, 일반 Bessel 적분 & 모두 유도과정에만 쓰며 해당 식 직전에 정의된다. 본문에
유도한 정확한 Gamma/Bessel 항등식으로 계산한다.\\
$\Gamma(\cdot),\zeta(\cdot),I_\mu,K_\mu$ & Gamma, Riemann zeta, 수정 Bessel 제1종과 제2종 &
검증된 special-function library를 사용하고 $K_\mu$는 식~\eqref{eq:K_definition}으로
교차검산한다.\\
$\Zsh_\nu,\Bmn_q,\Delta\Bmn_q$ & 이동 Epstein 합, 닫힌 지수합, registry 차분 & 식
\eqref{eq:Z_definition}, \eqref{eq:Bq}, \eqref{eq:DeltaBq}; cutoff를 늘려 수렴시키고
직접합과 비교한다.\\
$\beta,\gamma$ in 식~\eqref{eq:Bessel_generic_integral} & 그 적분에서만 쓰는 양의 변수 &
$\beta=\eta^2$, $\gamma=\pi^2\ell^2$이며 $\beta_T,\gamma_{m,n}$과 다른 기호다.\\
$\begin{gathered}\pi,e,\ii,\operatorname{Re}\\|\cdot|,\exp,\ln,\sum\\
\int,\lim,\max\end{gathered}$ & 원주율, 자연로그 밑, 허수단위, 실수부,
크기와 표준연산자 & $\exp x=e^x$, $\ln$은 그 역함수이며 각 합과 적분의 범위는 식에
표시한다.\\
$!,\bmod,\sim,\propto,\mathbb R$ & 팩토리얼, 나머지연산, 점근/동등표시, 비례표시,
실수집합 & 정수 $j\geq0$에서 $j!=\Gamma(j+1)$. $s\bmod b$는 셀 내부위치만 남기고
누적 장벽통과 이력은 unwrapped $s$ 또는 $z$에 보존한다.\\
\end{longtable}

\begin{longtable}{>{\raggedright\arraybackslash}p{0.17\textwidth}
                  >{\raggedright\arraybackslash}p{0.34\textwidth}
                  >{\raggedright\arraybackslash}p{0.39\textwidth}}
\toprule
동역학 기호 & 의미와 단위 & 정의와 계산방법\\
\midrule
\endfirsthead
\toprule
동역학 기호 & 의미와 단위 & 정의와 계산방법\\
\midrule
\endhead
\bottomrule
\endfoot
$t,f,\omega_{\rm L}$ & 시간, 반복주파수, 하중각주파수 & $\omega_{\rm L}=2\pi f$이며
비사인파는 의미 있는 모든 harmonic을 검사한다.\\
$T,\kB,\beta_T$ & 절대온도, Boltzmann 상수, 역열에너지 & $\beta_T=(\kB T)^{-1}$.\\
$q_i,v_i,m_i^*$ & 집단좌표, 속도, 유효질량 & $q_a=a,q_s=s$. 두 rigid patch가 함께
움직이면 $m_i^*=M_UM_L/(M_U+M_L)$, 한쪽 고정이면 이동쪽 질량을 쓴다.\\
$\mathcal K_i,R_i,\tau_{K_i}$ & memory-friction kernel, projected random force, kernel 감쇠시간 &
구속평형 MD에서 $\mathcal K_i(t)=\beta_T\langle R_i(0)R_i(t)\rangle$를 구하고 감쇠시간을
측정한다.\\
$\Gamma_i,M_i,D_i$ & Markov 마찰, 장시간 mobility, 확산계수 & Markov plateau가 있을 때
$\Gamma_i=\int_0^\infty\mathcal K_i\dd t$, $M_i=1/\Gamma_i$, $D_i=\kB TM_i$.\\
$\xi_i,\delta_{ij},\delta(t-t')$ & 단위 백색잡음, Kronecker delta, Dirac delta & 잡음상관은
$\langle\xi_i(t)\xi_j(t')\rangle=\delta_{ij}\delta(t-t')$.\\
$\rho,\rho_*,\mathcal P_b,\dd\Lambda$ & 위상공간밀도, 기준밀도, 결합 위상영역, 위상체적 &
$\mathcal P_b=\Omega_b\times\mathbb R^2$에서 정규화하며 $\rho_*$는 로그를 무차원화한다.\\
$\mathcal J_{v_i}^{\rm irr},\mathcal F_K,\dot D_K$ & 속도공간 비가역전류, 관성 자유에너지,
관성 소산율 & 식~\eqref{eq:velocity_irreversible_current},
\eqref{eq:Kramers_free_energy}, \eqref{eq:Kramers_dissipation}.\\
$\epsilon_i^{\rm in},\epsilon_i^{\rm mem}$ & 관성 및 memory 축약 판정수 & 식
\eqref{eq:overdamped_conditions}의 두 값이 모든 중요 하중주파수에서 작을 때만
Smoluchowski를 쓴다.\\
$P,P_{\rm eq},P_*,\mathcal Z_{\rm conf}$ & 상태밀도, 평형밀도, 기준밀도, 분배함수 &
$P=\iint\rho\dd v_a\dd v_s$; $P_{\rm eq}$는 유한 결합우물 안에서만 정규화한다.\\
$\begin{gathered}\Omega_b,\partial\Omega_b\\\partial\Omega_{\rm fr},\bm n_\Omega,\dd S\end{gathered}$ & 결합영역, 전체경계,
흡수형 균열경계, 외향법선, 경계면적요소 & 균열경계는 fitting한 $a_c$가 아니라 전체
에너지의 $a^\ddagger$에 둔다.\\
$J_a,J_s,\bm J$ & $a,s$방향 확률유속 & 식~\eqref{eq:Ja}--\eqref{eq:Js}; 보존형
finite-volume face flux로 계산한다.\\
$\mathcal F,\mu_P,D_{\rm irr},\dot D_{\rm irr}$ & 과감쇠 자유에너지, 확률 chemical
potential, 누적소산, 소산율 & 식~\eqref{eq:free_energy_functional},
\eqref{eq:probability_chemical_potential}, \eqref{eq:dissipation},
\eqref{eq:accumulated_dissipation}.\\
$S,P_{\rm crack},p_{\rm tail}$ & 생존확률, 누적균열확률, 순간 꼬리확률 & 식
\eqref{eq:survival_flux}, \eqref{eq:tail_probability}; 소산에너지를 균열확률과 동일시하지
않는다.\\
$\bar a,\bar s,\operatorname{Var}[a],\bar E$ & 평균거리, 평균 registry, 거리분산,
평균 potential 에너지 & 식~\eqref{eq:mean_a}--\eqref{eq:mean_energy}; 생존조건부 값은
$S$로 나눈다.\\
$g,\bm\nabla,\langle\cdot\rangle,\partial_t,\dot{(\cdot)},(\cdot)'$ & 일반 관측량, gradient,
ensemble 평균, 편미분, 전미분, 단일변수 미분 & 흡수경계에서는 식
\eqref{eq:general_moment_boundary}의 경계항을 유지한다.\\
\end{longtable}

\begin{longtable}{>{\raggedright\arraybackslash}p{0.17\textwidth}
                  >{\raggedright\arraybackslash}p{0.34\textwidth}
                  >{\raggedright\arraybackslash}p{0.39\textwidth}}
\toprule
소성/EAM 기호 & 의미와 단위 & 정의와 계산방법\\
\midrule
\endfirsthead
\toprule
소성/EAM 기호 & 의미와 단위 & 정의와 계산방법\\
\midrule
\endhead
\bottomrule
\endfoot
$s_z,s_{z,\pm}^\ddagger$ & $z$번째 안정점과 전진/후진 안장점; 길이 & 각 주기에서
$\partial_s\gamma=\tau$를 풀고 2차미분 부호로 분류한다.\\
$\gamma_p,\varepsilon_p$ & 축약 소성전단변형률과 축방향 투영 소성변형률 & 식
\eqref{eq:plastic_strain}; 전위경화식이 아니라 이상 단일슬립 bookkeeping이다.\\
$s_L,s_R,x,y,T_{\rm FP}$ & 반사경계, 흡수경계, 적분변수, 평균최초통과시간 & 정확한 우물과
안장점 주변 경계를 정하고 식~\eqref{eq:exact_MFPT}를 적응형 적분한다.\\
$k_{z,\pm}^{\rm FP},k_{z,\pm}$ & 최초통과 및 선택적 점근 전이율; 1/시간 & 과감쇠
MFPT의 역수 또는 높은 장벽에서만 식~\eqref{eq:Kramers_rate}.\\
$p_z,J_{a,z}$ & $z$번째 우물 확률밀도와 법선유속 & 시간척도 분리가 검증될 때만
식~\eqref{eq:master_equation}을 쓴다.\\
$\begin{gathered}E_{\rm EAM},\phi,F,\rho_i\\f,f_0,\beta_\rho\end{gathered}$ & EAM 에너지, pair항, embedding 함수,
전자밀도, density kernel, 진폭, 감쇠계수 & 완전하고 검증된 EAM parameter set을
사용하고 $f$를 먼저 합한 뒤 비선형 $F$를 계산한다.\\
$\mathcal U,\mathcal L,\mathcal U_{\rm cell}$ & 위와 아래 반결정 및 위쪽 interface cell &
선택한 결정방향에서 생성하고 이웃범위 또는 EAM cutoff 수렴성을 확인한다.\\
$M_U,M_L$ & 위와 아래 대표 patch 질량 & 상관영역 안 원자질량을 합산한다.\\
\end{longtable}
\normalsize

\subsection*{\koreanfont 실제 계산 순서}

\renewcommand{\labelenumi}{{\normalfont\arabic{enumi}.}}
\begin{enumerate}
  \item $m,n,\epsLJ,\sigLJ,b,A_{\rm at},A_0,s_0$와 normal-only/slip 분기를 선언한다.
  \item $U_\infty,W_{m,n}$을 계산하고 Bessel cutoff를 늘려 에너지와 필요한 미분값을
  수렴시킨 뒤 대표점에서 직접 원자합과 비교한다.
  \item 전체 에너지의 근을 bracket하여 안정점, 안장점, 불안정점을 구하고 정확한
  곡률로 분류한다. 에너지를 다항식으로 fitting할 필요가 없다.
  \item 원자기하와 구속평형 MD에서 $m_i^*,\mathcal K_i$를 구하고
  식~\eqref{eq:overdamped_conditions}을 검사한다. 조건이 맞지 않으면 mobility를
  기본상수처럼 넣지 않는다.
  \item 관성이 중요하면 Kramers PDE를 푼다. 축약이 검증되면 Smoluchowski PDE를
  쓴다. 두 경우 모두 보존형 finite-volume 기법을 적용하고, 격자, 시간간격,
  속도영역, 정규화의 수렴을 확인한다.
  \item 균열경계는 $a^\ddagger$에 두고 바깥 확률유속을 적분한다. $p_{\rm tail}$,
  $S$, $P_{\rm crack}$, $D_{\rm irr}$를 서로 다른 출력으로 보고한다.
\end{enumerate}


\section*{\koreanfont 초록}

이 문서는 법선방향 거리 $a$만 사용하는 기존 피로모델에 상대 미끄럼 좌표 $s$를
추가하고, 일반화된 Lennard--Jones $m$--$n$ 상호작용으로부터 정확한 슬립 의존
격자에너지 $W_{m,n}(a,s)$를 유도한다. 두 무한 원자열의 직접합은 이동된
Epstein--Hurwitz 합으로 표현되며, Mellin 변환과 Poisson 합공식을 거치면
수정 Bessel 제2종 함수 $K_\mu$를 포함하는 지수수렴 급수로 정확히 변환된다.
이 과정에서는 $a$나 $s$에 대한 Taylor 전개를 사용하지 않는다. 이어서 외력까지
포함한 전체 유효에너지, Smoluchowski 확률수송식, 평균거리와 평균에너지,
비가역 에너지 축적식, 정확한 최초통과시간, 소성 슬립의 조건을 정리한다.

\section*{\koreanfont 번역 1. 변수, 에너지 계수법, 적용범위}

$a>0$는 원자 또는 원자층의 법선방향 거리, $s$는 지정된 미끄럼방향의 상대 변위,
$b$는 그 방향의 격자 반복거리이다. 두 지수는
\[
  m>n>1
\]
을 만족해야 한다. $m>n$은 짧은 거리에서 반발항이 인력항보다 강해야 한다는 조건이고,
$n>1$은 일차원 무한급수 $\sum k^{-n}$이 수렴하기 위한 조건이다.

유한 원자계의 쌍에너지는 각 쌍을 한 번씩만 세어야 하므로
\[
  E_{\rm pair}
  =\frac12\sum_i\sum_{j\ne i}v(r_{ij})
  =\sum_{i<j}v(r_{ij})
\]
이다. 모든 원자에서 양방향 이웃을 더할 때에는 앞의 $1/2$가 반드시 필요하다.
반면 하나의 위쪽 기준원자와 아래쪽 원자열 사이의 교차쌍만 더하는 $W(a,s)$에서는
각 교차쌍이 애초에 한 번만 등장하므로 $1/2$를 붙이지 않는다.

$s$는 유한요소 좌표의 차원을 늘리는 변수가 아니라 재료점 내부의 미시상태변수이다.
따라서 구조물 메시는 1D, 2D, 3D로 만들 수 있지만, 현재 구성법칙은 선택된 법선응력과
필요할 때만 하나의 resolved slip 좌표를 사용한다. 엄격한 법선피로 분기에서는
$s=s_0$, $\tau=0$으로 두면 된다.

\section*{\koreanfont 번역 2. 일반화된 Lennard--Jones $m$--$n$ potential}

$\epsLJ$를 우물깊이, $\sigLJ$를 에너지가 영이 되는 거리로 유지하려면 일반형을
\[
  \boxed{
  v_{m,n}(r)=C_{m,n}\epsLJ
  \left[
    \left(\frac{\sigLJ}{r}\right)^m
    -\left(\frac{\sigLJ}{r}\right)^n
  \right],\qquad
  C_{m,n}=\frac{m}{m-n}
  \left(\frac{m}{n}\right)^{n/(m-n)}}
\]
로 정규화해야 한다. $v'(r_e)=0$을 풀면
\[
  r_e=\sigLJ\left(\frac{m}{n}\right)^{1/(m-n)}
\]
이고, 이 값을 $v(r_e)=-\epsLJ$에 대입하면 위의 $C_{m,n}$이 나온다. 같은 potential을
평형거리로 쓰면
\[
  v_{m,n}(r)=\epsLJ\left[
  \frac{n}{m-n}\left(\frac{r_e}{r}\right)^m
  -\frac{m}{m-n}\left(\frac{r_e}{r}\right)^n
  \right]
\]
이다. 따라서 지수에 숫자를 넣지 않아도 $v(r_e)=-\epsLJ$, $v'(r_e)=0$,
$v''(r_e)>0$이 정확히 성립한다.

\subsection*{\koreanfont 양방향 무한 원자열의 실제 에너지}

등간격 $a$를 갖는 $N$개 원자열의 전체에너지는
\[
  E_N(a)=\sum_{k=1}^{N-1}(N-k)v_{m,n}(ka)
\]
이다. 거리 $ka$만큼 떨어진 쌍이 정확히 $N-k$개이기 때문이다. 이를 $N$으로 나누고
$N\to\infty$를 취하면 원자당 에너지는
\[
  \boxed{
  U_\infty(a)=C_{m,n}\epsLJ
  \left[
    \left(\frac{\sigLJ}{a}\right)^m\zeta(m)
    -\left(\frac{\sigLJ}{a}\right)^n\zeta(n)
  \right]}
\]
이 된다. 내부 기준원자 하나가 양쪽 이웃과 상호작용하는 국소합은 $2U_\infty$지만,
전체에너지에서는 모든 쌍이 두 번 세어졌으므로 $1/2$를 곱해 다시 $U_\infty$가 된다.
즉 위 식은 양방향을 누락한 식이 아니라 중복계수까지 올바르게 처리한 원자당
총에너지이다.

무하중 격자평형 $U_\infty'(a_0)=0$으로부터
\[
  \boxed{
  a_0=\sigLJ
  \left[\frac{m\zeta(m)}{n\zeta(n)}\right]^{1/(m-n)}}
\]
을 얻는다. 고립된 두 원자의 $r_e$와 무한격자의 $a_0$가 다른 이유는 후자에는 모든
원거리 이웃의 상호작용이 포함되기 때문이다.

\section*{\koreanfont 번역 3. 두 원자열의 기하와 직접 격자합}

미끄럼방향 단위벡터를 $\bm d$, 법선벡터를 $\bm n$이라 하고
$\bm d\cdot\bm n=0$으로 둔다. 아래쪽 원자의 위치가 $pb\bm d$일 때 위쪽
기준원자와 $p$번째 아래 원자 사이 거리는
\[
  \boxed{r_p(a,s)=\sqrt{a^2+(pb+s)^2}}
\]
이다. 따라서 교차 원자열 에너지는 근사 없이
\[
  W_{m,n}(a,s)=C_{m,n}\epsLJ
  \sum_{p=-\infty}^{\infty}\left[
  \frac{\sigLJ^m}{[a^2+(pb+s)^2]^{m/2}}
  -\frac{\sigLJ^n}{[a^2+(pb+s)^2]^{n/2}}
  \right]
\]
로 주어진다. 큰 $|p|$에서 각 항은 각각 $|p|^{-m}$, $|p|^{-n}$처럼 감소하므로
$m>n>1$이면 두 급수가 각각 절대수렴한다.

무차원 변수
\[
  \delta=\frac{s}{b},\qquad \eta=\frac{a}{b}
\]
와 이동된 합
\[
  \boxed{
  \Zsh_\nu(\delta,\eta)=
  \sum_{p=-\infty}^{\infty}
  [(p+\delta)^2+\eta^2]^{-\nu}}
\]
를 정의하면
\[
  \boxed{
  W_{m,n}(a,s)=C_{m,n}\epsLJ
  \left[
  \left(\frac{\sigLJ}{b}\right)^m\Zsh_{m/2}
  -\left(\frac{\sigLJ}{b}\right)^n\Zsh_{n/2}
  \right]}
\]
이다. $p$의 지표를 하나 이동해도 무한합이 같으므로
$W(a,s+b)=W(a,s)$가 정확히 성립한다. 이 주기성이 서로 다른 격자 registry 사이의
에너지 장벽을 만든다.

\section*{\koreanfont 번역 4. Mellin--Poisson--Bessel 변환의 상세 유도}

\subsection*{\koreanfont 4.1 Mellin 표현}

Gamma 함수의 정의
\[
  \Gamma(\nu)=\int_0^\infty u^{\nu-1}e^{-u}\,\dd u
\]
에서 $u=xt$를 대입하면
\[
  x^{-\nu}=\frac1{\Gamma(\nu)}
  \int_0^\infty t^{\nu-1}e^{-xt}\,\dd t
\]
를 얻는다. 즉 역거듭제곱을 Gaussian들의 연속중첩으로 정확히 바꾼 것이다.
이를 $x=(p+\delta)^2+\eta^2$에 적용하면
\[
  \Zsh_\nu=
  \frac1{\Gamma(\nu)}\int_0^\infty
  t^{\nu-1}e^{-\eta^2t}
  \sum_{p\in\mathbb Z}e^{-t(p+\delta)^2}\,\dd t .
\]
실수 $\nu>1/2$에서는 피적분함수가 음수가 아니므로 Tonelli 정리에 의해 합과 적분을
바꿀 수 있다. 복소수 $\nu$에서도 같은 반평면의 절대수렴성이 Fubini 정리를 보장한다.

\subsection*{\koreanfont 4.2 이동 Gaussian과 Poisson 합공식}

$g(x)=e^{-t(x+\delta)^2}$의 Fourier 변환을
$\widehat g(\ell)=\int g(x)e^{-2\pi\ii\ell x}\dd x$로 정의한다.
$y=x+\delta$를 넣고 제곱을 완성하면
\[
  -ty^2-2\pi\ii\ell y
  =-t\left(y+\frac{\pi\ii\ell}{t}\right)^2
   -\frac{\pi^2\ell^2}{t}
\]
이므로
\[
  \widehat g(\ell)=
  \sqrt{\frac\pi t}\,
  e^{-\pi^2\ell^2/t}e^{2\pi\ii\ell\delta}.
\]
따라서 Poisson 합공식은
\[
  \boxed{
  \sum_{p\in\mathbb Z}e^{-t(p+\delta)^2}
  =\sqrt{\frac\pi t}
  \sum_{\ell\in\mathbb Z}
  e^{-\pi^2\ell^2/t}e^{2\pi\ii\ell\delta}}
\]
가 된다. 이를 앞의 Mellin 적분에 넣으면 원자공간의 지표 $p$가 역격자 지표
$\ell$로 바뀐다.

\subsection*{\koreanfont 4.3 수정 Bessel 제2종 함수 $K_\mu$의 정의}

$K_\mu(\chi)$는 다음 수정 Bessel 방정식
\[
  \chi^2y''+\chi y'-(\chi^2+\mu^2)y=0
\]
의 두 독립해 중 큰 $\chi$에서 감쇠하는 해이다. Bessel 인자는 뒤의 우물번호 $z$와
혼동하지 않도록 $\chi$로 표기한다. 이 문서에서 사용하는 직접적인 정의는
\[
  \boxed{
  K_\mu(\chi)=\frac12\int_0^\infty
  u^{\mu-1}
  \exp\left[-\frac \chi2\left(u+\frac1u\right)\right]\dd u,
  \qquad \chi>0}
\]
이다. $u\mapsto1/u$를 대입하면 $K_{-\mu}(\chi)=K_\mu(\chi)$를 바로 확인할 수 있다.
실수 $\mu$와 양의 $\chi$에서는 피적분함수가 양수이므로 $K_\mu(\chi)>0$이다.
또한 큰 $\chi$에서
\[
  K_\mu(\chi)\sim\sqrt{\frac{\pi}{2\chi}}e^{-\chi}
\]
이므로 뒤에 나오는 역격자 급수가 지수적으로 빠르게 수렴한다.

이제
\[
  \mathcal I=\int_0^\infty
  t^{\mu-1}e^{-\beta t-\gamma/t}\dd t
\]
에서 $t=\sqrt{\gamma/\beta}\,u$를 대입한다. 그러면
\[
  t^{\mu-1}\dd t=
  \left(\frac\gamma\beta\right)^{\mu/2}
  u^{\mu-1}\dd u,
  \qquad
  \beta t+\frac\gamma t=
  \sqrt{\beta\gamma}\left(u+\frac1u\right).
\]
따라서 위 적분은 $K_\mu$의 정의와 항별로 정확히 일치하여
\[
  \boxed{
  \mathcal I=
  2\left(\frac\gamma\beta\right)^{\mu/2}
  K_\mu(2\sqrt{\beta\gamma})}
\]
가 된다. 이것은 표에서 가져온 경험식도 아니고 Taylor 근사도 아니다.

현재 문제에서는
\[
  \mu=\nu-\frac12,\qquad
  \beta=\eta^2,\qquad
  \gamma=\pi^2\ell^2
\]
이므로
\[
  I_\ell=
  2\left(\frac{\pi|\ell|}{\eta}\right)^{\nu-1/2}
  K_{\nu-1/2}(2\pi|\ell|\eta).
\]
$\ell$과 $-\ell$을 합치면 복소지수는 cosine으로 바뀌며 최종적으로
\[
  \boxed{
  \Zsh_\nu(\delta,\eta)=
  \frac{\sqrt\pi}{\Gamma(\nu)}
  \left[
  \Gamma\left(\nu-\frac12\right)\eta^{1-2\nu}
  +4\sum_{\ell=1}^{\infty}
  \cos(2\pi\ell\delta)
  \left(\frac{\pi\ell}{\eta}\right)^{\nu-1/2}
  K_{\nu-1/2}(2\pi\ell\eta)
  \right]}
\]
을 얻는다. 첫 항은 registry와 무관한 영 Fourier mode이고, 뒤의 급수만 $s$에 따른
주기적 에너지 장벽을 만든다.

\section*{\koreanfont 번역 5. 일반 $m,n$에 대한 완전한 에너지식}

$q>1$에 대해
\[
  \boxed{
  \begin{aligned}
  \Bmn_q(\delta,\eta)=
  \frac{\sqrt\pi}{\Gamma(q/2)}
  \Bigg[
  &\Gamma\left(\frac{q-1}{2}\right)\eta^{1-q}\\
  &+4\sum_{\ell=1}^{\infty}
  \cos(2\pi\ell\delta)
  \left(\frac{\pi\ell}{\eta}\right)^{(q-1)/2}
  K_{(q-1)/2}(2\pi\ell\eta)
  \Bigg]
  \end{aligned}}
\]
로 두면 $\Bmn_q=\Zsh_{q/2}$이다. 따라서 완전한 두 원자열 에너지는
\[
  \boxed{
  W_{m,n}(a,s)=C_{m,n}\epsLJ
  \left[
  \left(\frac{\sigLJ}{b}\right)^m\Bmn_m(\delta,\eta)
  -\left(\frac{\sigLJ}{b}\right)^n\Bmn_n(\delta,\eta)
  \right]}
\]
이다. $\Bmn_m$과 $\Bmn_n$의 모든 항을 한 식에 펼친 형태가 원문
식~\eqref{eq:full_W_mn}이다. $\sigLJ/b$, $\delta$, $\eta$는 모두 무차원이므로
$W$의 단위는 $\epsLJ$와 같은 에너지이다.

\section*{\koreanfont 번역 6. 초과 슬립에너지, 내부전단응력, 정확한 장벽}

완전한 registry를 $s_0$라 하면 초과에너지는
\[
  \Delta W_{m,n}(a,s)=W_{m,n}(a,s)-W_{m,n}(a,s_0)
\]
이고 원자 하나가 대표하는 미끄럼면 면적을 $A_{\rm at}$라 할 때
\[
  \boxed{
  \gamma_{m,n}(a,s)=\frac{\Delta W_{m,n}(a,s)}{A_{\rm at}}}
\]
이다. 차분을 취하면 registry와 무관한 영 Fourier mode가 정확히 사라진다.
따라서 슬립 장벽은 전적으로 cosine 급수에서 나온다.

내부 저항전단응력은 에너지면적밀도를 변위로 미분하여
\[
  \boxed{
  \tau_{\rm int}(a,s)=
  \frac{\partial\gamma_{m,n}}{\partial s}
  =\frac{C_{m,n}\epsLJ}{A_{\rm at}b}
  \left[
  \left(\frac{\sigLJ}{b}\right)^m\partial_\delta\Bmn_m
  -\left(\frac{\sigLJ}{b}\right)^n\partial_\delta\Bmn_n
  \right]}
\]
이다. 여기서 $\partial_s=b^{-1}\partial_\delta$이며
\[
  \partial_\delta\Bmn_q=
  -\frac{8\pi\sqrt\pi}{\Gamma(q/2)}
  \sum_{\ell=1}^\infty
  \ell\sin(2\pi\ell\delta)
  \left(\frac{\pi\ell}{\eta}\right)^{(q-1)/2}
  K_{(q-1)/2}(2\pi\ell\eta).
\]
단위는 $(\text{에너지}/\text{면적})/\text{길이}=\text{응력}$으로 일치한다.

면적 $A_0$인 영역에 resolved shear $\tau$가 작용하면 기울어진 에너지는
\[
  \psi(a,s,t)=A_0\gamma_{m,n}(a,s)-\tau(t)A_0(s-s_0)
\]
이다. 안정점은 $\partial_s\gamma=\tau$, $\partial_s^2\gamma>0$을 만족하고,
안장점은 같은 힘평형과 $\partial_s^2\gamma<0$을 만족한다. 따라서 정확한 전이장벽은
\[
  \boxed{
  \Delta G_{z,\pm}^\ddagger
  =A_0\left[
  \gamma(a,s_{z,\pm}^\ddagger)-\gamma(a,s_z)
  -\tau(s_{z,\pm}^\ddagger-s_z)
  \right]}
\]
이다. 국소 최소와 안장점이 합쳐져
$\partial_s\gamma=\tau$, $\partial_s^2\gamma=0$이 되는 순간이 무열적
격자불안정 조건이다.

\section*{\koreanfont 번역 7. 외력을 포함한 대표영역 전체에너지}

면적 $A_0$인 대표영역에 $N_A=A_0/A_{\rm at}$개의 등가 원자기둥이 있다고 하자.
법선 기준에너지에는 $N_AU_\infty$를 쓰고, 슬립부분에는 기준 registry에 대한
초과에너지만 더하면 중복계산을 피할 수 있다. 전체 축약 potential은
\[
  \boxed{
  \begin{aligned}
  \Phi(a,s,t)={}&N_A C_{m,n}\epsLJ\left[
  \left(\frac{\sigLJ}{a}\right)^m\zeta(m)
  -\left(\frac{\sigLJ}{a}\right)^n\zeta(n)
  \right]\\
  &+A_0\gamma_{m,n}(a,s)
  -\sigma(t)A_0(a-a_0)
  -\tau(t)A_0(s-s_0).
  \end{aligned}}
\]
외부응력을 원자힘으로 바꾸는 면적은 임의계수가 아니다.
$F_n=\sigma A_0$, $F_s=\tau A_0$가 traction의 정의에서 바로 나온다.
단축응력에서 특정 slip system의 resolved shear는
$\tau=\mathcal M\sigma$로 쓸 수 있다.

고정된 하중과 유한한 결합영역 $\Omega_b$ 안에서는
\[
  P_{\rm eq}(a,s)=\frac1{\mathcal Z_{\rm conf}}
  e^{-\beta_T\Phi(a,s)},\qquad
  \beta_T=(\kB T)^{-1}
\]
이다. 그러나 dead tensile stress에서는 $-\sigma A_0a$ 때문에
$a\to\infty$에서 $\Phi\to-\infty$가 되어 전 구간 평형분포를 정규화할 수 없다.
그러므로 균열 전 상태는 결합우물 안의 준평형분포와 흡수경계로 다뤄야 한다.

외부 장벽 $a^\ddagger$는 임의의 임계거리로 두지 않고
\[
  N_AU_\infty'(a^\ddagger)=\sigma A_0,\qquad
  U_\infty''(a^\ddagger)<0
\]
으로 정한다. 안정점과 장벽이 합쳐지는 격자 변곡점은
\[
  \boxed{
  a_{\rm inf}=\sigLJ
  \left[
  \frac{m(m+1)\zeta(m)}
       {n(n+1)\zeta(n)}
  \right]^{1/(m-n)}}
\]
이다. 따라서 물리적으로 유도된 꼬리확률은
\[
  p_{\rm tail}(t;\sigma)=
  \int_{a^\ddagger(\sigma)}^\infty\int P(a,s,t)\dd s\dd a
\]
이며, 흡수경계를 사용하면 그 경계를 통과한 누적 확률유속이 균열개시확률이 된다.

\section*{\koreanfont 번역 8. 관성 확률동역학, 축약조건, 히스테리시스}

원자 수준에서 힘은 속도가 아니라 가속도를 만든다. 따라서 이 문서의 기본
coarse-grained 식은 mobility 식이 아니라 memory와 관성을 포함한
\[
  \boxed{
  m_i^*\ddot q_i(t)=-\partial_{q_i}\Phi
  -\int_0^t\mathcal K_i(t-t')\dot q_i(t')\dd t'+R_i(t)}
\]
이다. 여기서 $q_a=a$, $q_s=s$, $m_i^*$는 집단좌표의 유효질량,
$\mathcal K_i$는 제거된 원자진동 및 phonon 자유도의 memory-friction kernel이고,
$R_i$는 그 자유도에서 되돌아오는 무작위 힘이다. 열평형에서는
\[
  \langle R_i(t)R_j(t')\rangle
  =\kB T\mathcal K_i(|t-t'|)\delta_{ij}
\]
가 성립한다.

memory가 충분히 짧을 때에만 마찰을 $\Gamma_i\dot q_i$로 바꾸어
\[
  \dot q_i=v_i,\qquad
  m_i^*\dot v_i=-\partial_{q_i}\Phi-\Gamma_iv_i
  +\sqrt{2\Gamma_i\kB T}\,\xi_i
\]
라는 Markov Langevin 식을 얻는다. 이에 대응하는 기본 확률밀도는
$P(a,s,t)$가 아니라 위치와 속도를 모두 포함한
$\rho(a,s,v_a,v_s,t)$이며, 그 지배식은 영문 식~\eqref{eq:Kramers_PDE}의
Kramers 방정식이다. 우리가 원하는 거리--슬립 분포는
\[
  \boxed{P(a,s,t)=\iint\rho(a,s,v_a,v_s,t)\dd v_a\dd v_s}
\]
로 얻는다. 일반적으로 속도를 적분한 것만으로는 닫힌 PDE가 바로 생기지 않는다.

관성 Kramers 모델의 자유에너지는 운동에너지까지 포함한다.
\[
  \mathcal F_K=\int\rho\left[
  \Phi+\frac12m_a^*v_a^2+\frac12m_s^*v_s^2
  +\kB T\ln\left(\frac{\rho}{\rho_*}\right)
  \right]\dd\Lambda .
\]
속도공간의 비가역 전류 $\mathcal J_{v_i}^{\rm irr}$를 사용하면
\[
  \boxed{
  \frac{\dd\mathcal F_K}{\dd t}
  =\langle\partial_t\Phi\rangle_\rho-\dot D_K,qquad
  \dot D_K=\sum_i\int
  \frac{(m_i^*)^2}{\Gamma_i\rho}
  (\mathcal J_{v_i}^{\rm irr})^2\dd\Lambda\ge0}
\]
이다. 이것이 관성을 유지한 Markov 모델의 소산식이다.

Smoluchowski 식을 쓰려면 먼저
\[
  \boxed{
  \epsilon_i^{\rm in}=\frac{(2\pi f)m_i^*}{\Gamma_i}\ll1,
  \qquad
  \epsilon_i^{\rm mem}=(2\pi f)\tau_{K_i}\ll1}
\]
을 확인해야 한다. 이 조건에서만
\[
  M_i=\Gamma_i^{-1},\qquad D_i=\kB TM_i
\]
를 정의하고 관성을 제거하여
\[
  \boxed{\partial_tP=-\partial_aJ_a-\partial_sJ_s}
\]
\[
  J_a=-M_a(P\partial_a\Phi+\kB T\partial_aP),\qquad
  J_s=-M_s(P\partial_s\Phi+\kB T\partial_sP)
\]
를 사용한다. 따라서 $\dot q=-M\partial_q\Phi$는 원자의 기본법칙이 아니라
빠른 관성운동이 제거된 장시간 terminal-response 식이다.

닫힌 결합영역에서는 $\iint P\dd a\dd s=1$이다. 흡수형 균열경계에서는
\[
  S(t)=\iint_{\Omega_b}P\dd a\dd s,qquad
  P_{\rm crack}(t)=1-S(t)
\]
이며 $P_{\rm crack}$은 흡수경계의 바깥방향 확률유속을 시간적분해서 계산한다.
평균거리와 평균에너지는
\[
  \bar a=\iint aP\dd a\dd s,qquad
  \bar E=\iint\Phi P\dd a\dd s
\]
이고, 흡수문제에서 생존한 상태만의 조건부 평균을 원하면 각각 $S(t)$로 나눈다.
비선형 $\Phi$를 이차 Taylor 전개로 닫으면 꼬리와 장벽을 잃으므로 $P$ 또는 $\rho$를
직접 전진시키는 것이 원칙이다.

과감쇠 자유에너지는
\[
  \mathcal F[P,t]=\iint
  \left[\Phi P+\kB TP\ln\left(\frac{P}{P_*}\right)\right]\dd a\dd s
\]
이다. no-flux 또는 periodic 경계에서만 기존 식
\[
  \boxed{
  \frac{\dd\mathcal F}{\dd t}
  =\langle\partial_t\Phi\rangle-\dot D_{\rm irr},\qquad
  \dot D_{\rm irr}=\iint\left[
  \frac{J_a^2}{M_aP}+\frac{J_s^2}{M_sP}
  \right]\dd a\dd s\ge0}
\]
이 그대로 성립한다. 흡수형 균열경계에서는 반드시
\[
  -\int_{\partial\Omega_b}\mu_P\bm J\mathbin{\cdot}\bm n_\Omega\dd S
\]
라는 자유에너지 경계유출항을 추가해야 한다. $P\to0$ 근처에서는 나눗셈 형태보다
\[
  \dot D_{\rm irr}=\iint\sum_iM_iP(\partial_{q_i}\mu_P)^2\dd a\dd s
\]
를 사용하는 편이 안전하다.

따라서 loading과 unloading의 동적 차이는 유한한 관성, 마찰, 확률완화에서 나오고,
영구 소성변형은 추가로 우물 사이 $z$ 전이를 요구한다. 순수 normal-only 분기에서는
$s=s_0$, $J_s=0$으로 두고 $(a,v_a)$ Kramers 식 또는 조건이 검증된 $a$방향
Smoluchowski 식만 사용한다.

\section*{\koreanfont 번역 9. 이 모델로 소성변형을 표현할 수 있는가?}

\subsection*{\koreanfont 9.1 격자주기를 풀어 쓴 소성 슬립}

$s$를 항상 $b$로 나눈 나머지만 저장하면 몇 번의 격자장벽을 넘었는지 알 수 없다.
따라서
\[
  \boxed{
  s(t)=z(t)b+\widetilde s(t),\qquad
  z\in\mathbb Z,\quad0\le\widetilde s<b}
\]
로 풀어 써야 한다. $\widetilde s$의 우물 내부 운동은 주로 회복가능하고, 정수
$z$의 변화가 비가역 slip event이다. 슬립이 $h_{\rm slip}$ 두께에 걸쳐
균질화된다면
\[
  \boxed{
  \gamma_p=\frac{b}{h_{\rm slip}}\langle z\rangle,
  \qquad
  \varepsilon_p=\mathcal M\gamma_p}
\]
이다. 따라서 unloading 뒤 $\widetilde s$가 원위치로 돌아와도 $z$가 남으면 영구변형이
존재한다.

\subsection*{\koreanfont 9.2 과감쇠 조건에서 Taylor 근사 없는 정확한 최초통과시간}

고정된 $a$와 frozen load에서 왼쪽 경계를 반사경계 $s_L$, 오른쪽 장벽 너머를
흡수경계 $s_R$로 둔다. 평균 최초통과시간 $T_{\rm FP}(s)$는 backward
Smoluchowski 식
\[
  -1=-M_s\psi'(s)T_{\rm FP}'(s)
  +M_s\kB TT_{\rm FP}''(s),
  \qquad
  T_{\rm FP}'(s_L)=0,\quad T_{\rm FP}(s_R)=0
\]
을 만족한다. 이 식은 식~\eqref{eq:overdamped_conditions}가 검증된 경우에만 적용한다.
여기에 $e^{-\beta_T\psi}$를 곱하면 왼쪽이 완전미분이 되며,
반사경계에서 한 번, 흡수경계까지 다시 한 번 적분하면
\[
  \boxed{
  T_{\rm FP}(s)=\frac1{M_s\kB T}
  \int_s^{s_R}e^{\beta_T\psi(y)}
  \left[
  \int_{s_L}^{y}e^{-\beta_T\psi(x)}\dd x
  \right]\dd y}
\]
를 얻는다. 이 식은 명시된 1D frozen-load 과감쇠 경계값 문제에서는 정확하고 full
$\psi$를 그대로 적분하므로 장벽을 포물선으로 근사하지 않는다. 그러나 관성 Kramers
동역학의 정확한 최초통과식은 아니다.
고정하중 유효전이율은 $k_{z,+}^{\rm FP}=1/T_{\rm FP}(s_z)$로 정의할 수 있다.
다만 이를 시간국소 master equation에 넣으려면 우물내 평형, 우물간 탈출, 하중변화
시간척도가 분리되어야 한다. 이 조건이 맞지 않으면 전이율 kernel로 줄이지 말고
시간의존 Smoluchowski PDE를 직접 풀어야 하며 관성이 중요하면 Kramers PDE를 푼다.

Kramers 식은 높은 장벽에서 위 정확식을 단순화한 선택적 점근식이다. 장벽높이는
full energy에서 계산하지만 prefactor는 최소점과 안장점 부근의 국소 조화근사를
사용한다. 여기서 Kramers 점근식은 전이율 공식을 뜻하며, 본 이론의 기본
Kramers PDE와 구분해야 한다. 기본 지배식은 관성 Kramers PDE이고, 검증된 과감쇠
극한에서만 Smoluchowski PDE와 위 최초통과 적분식을 사용한다.

well 확률 $p_z(a,t)$를 쓰는 축약 master equation은
\[
  \partial_tp_z=-\partial_aJ_{a,z}
  +k_{z-1,+}p_{z-1}+k_{z+1,-}p_{z+1}
  -(k_{z,+}+k_{z,-})p_z
\]
이다. 여기서
\[
  \frac{\dd\langle z\rangle}{\dd t}
  =\sum_z\int(k_{z,+}-k_{z,-})p_z\dd a
\]
이므로 $\langle z\rangle$의 변화가 곧 소성변형률의 변화가 된다.

\subsection*{\koreanfont 9.3 가능한 것과 아직 불가능한 것}

두 원자열 중앙력 모델은 주기적 lattice resistance, 열활성 registry jump,
이상적인 단일슬립 영구변형, 그리고 소산 히스테리시스를 표현할 수 있다.
그러나 전위선의 형상과 장거리 탄성장, 전위증식, forest hardening, backstress,
전위 junction, 입계, 여러 FCC slip system은 아직 포함하지 않는다.

따라서 $s$를 끄고 $a$만 쓰는 현재의 순수 법선피로 모델은 탄성개구, 준안정
decohesion, 분포꼬리 변화에는 사용할 수 있지만 결정학적 소성이라고 부를 수 없다.
$s$와 $z$를 켜면 이상적인 단일슬립 소성모델이 된다. 실제 알루미늄의 정량적
cyclic plasticity로 확장하려면 검증된 결정에너지면과 전위저장 및 경화 내부변수가
추가되어야 한다.

\section*{\koreanfont 번역 10. EAM으로 교체하는 경우}

금속결합을 더 현실적으로 다루려면 pair LJ를 embedded-atom method로 바꿀 수 있다.
EAM 전체에너지는
\[
  \boxed{
  E_{\rm EAM}(a,s)=
  \frac12\sum_i\sum_{j\ne i}\phi(r_{ij}(a,s))
  +\sum_iF(\rho_i(a,s)),\qquad
  \rho_i=\sum_{j\ne i}f(r_{ij})}
\]
이다. 제안된 전자밀도 kernel은
\[
  \boxed{
  f(r)=f_0\exp\left[
  -\beta_\rho\left(\frac r{r_e}-1\right)
  \right]}
\]
로 둘 수 있다. 여기서 $\beta_\rho$는 무차원 밀도감쇠계수이며 열역학의
$\beta_T$와 다른 변수이다.

계산할 때에는 먼저 각 이웃의 $f(r_{ij})$를 합하여 $\rho_i$를 만든다.

그 다음 이 합에 비선형 함수 $F(\rho_i)$를 적용한다. $f(r)$를
pair potential처럼 $F$ 안에 하나씩 대입하여 더하면 EAM의 many-body 성질이
사라지므로 잘못이다.

EAM generalized stacking-fault energy는
\[
  \boxed{
  \gamma_{\rm EAM}(a,s)=
  \frac{E_{\rm EAM}(a,s)-E_{\rm EAM}(a,s_0)}
       {A_{\rm cell}}}
\]
이다. 기존 확률수송식에서는 $\gamma_{m,n}$만 $\gamma_{\rm EAM}$으로 교체하면
구조가 그대로 유지된다. 다만 지수형 밀도합과 비선형 $F$는 일반적으로 동일한
Epstein $m,n$ 급수로 닫히지 않으므로 실제 neighbor geometry에서 직접 계산해야 한다.

\section*{\koreanfont 번역 11. FCC 결정 반공간으로의 확장}

선택된 면을 기준으로 위쪽 반결정 $\mathcal U$와 아래쪽 반결정 $\mathcal L$을
나누면 pair potential의 정확한 초과에너지면은
\[
  \boxed{
  \begin{aligned}
  \gamma_{\rm GSF}(a,s)=\frac1{A_{\rm cell}}
  \sum_{i\in\mathcal U_{\rm cell}}
  \sum_{j\in\mathcal L}
  \Big[
  &v_{m,n}(|\bm R_{ij}+a\bm n+s\bm d|)\\
  &-v_{m,n}(|\bm R_{ij}+a\bm n+s_0\bm d|)
  \Big].
  \end{aligned}}
\]
모든 cross-plane pair를 한 번씩만 세므로 여기에는 $1/2$가 없다. 실제 FCC에서는
하나의 원자열 지표 대신 두 개의 면내 격자지표와 여러 법선층이 들어간다.
역거듭제곱 pair항에는 같은 Mellin--Poisson 절차를 적용할 수 있지만 역격자합이
2차원이 된다. 정량적 알루미늄 모델에서는 선택된 FCC slip plane과 direction에 대한
EAM 또는 first-principles GSF를 사용하는 것이 타당하다.

구조물의 2D 또는 3D meshing 자체는 이 미시모델의 차원과 별개이다. 각 유한요소에서
선택한 법선응력 $\sigma_n(\bm x,t)$를 $\Phi$에 전달하면 된다. 전단소성을 사용하지
않는 분기에서는 $s$와 $\tau$를 비활성화한 채 3D CAD와 3D mesh를 그대로 사용할 수 있다.

\section*{\koreanfont 번역 12. 해석적 검산조건}

수치 대입 없이 다음 항등식을 확인할 수 있다.
\[
  \Bmn_q(\delta+1,\eta)=\Bmn_q(\delta,\eta),\qquad
  \Bmn_q(-\delta,\eta)=\Bmn_q(\delta,\eta),
\]
\[
  \partial_\delta\Bmn_q(0,\eta)=0,\qquad
  \Delta\Bmn_q(\delta_0,\delta_0;\eta)=0.
\]
직접 원자합과 Poisson--Bessel 급수는 모든 $q>1$, $\eta>0$, 실수 $\delta$에서
같아야 한다. $\eta$가 커지면 registry 의존항은 $K$ 함수 때문에 지수적으로 사라진다.
또한 법선방향 Riemann-zeta 원자열과 두 원자열 slip geometry는 서로 다른 기하이다.
$\eta\to0$, $\delta\to0$으로 법선 원자열을 얻으려 하면 두 원자가 겹쳐 발산하므로
그런 극한을 사용해서는 안 된다.

\section*{\koreanfont 번역 13. 연구식의 최종 정리}

현재 이론의 핵심식을 연구자가 사용해 온 네 식과 확률꼬리식에 맞춰 정리하면 다음과 같다.

\begingroup
\setlength{\itemsep}{0.1\baselineskip}
\setlength{\parsep}{0pt}
\setlength{\topsep}{0.4\baselineskip}
\setlength{\abovedisplayskip}{4pt}
\setlength{\belowdisplayskip}{4pt}
\par\smallskip
\begin{enumerate}
  \item \textbf{거리--에너지식:}
  $U_\infty(a)$가 순수 법선 원자열 에너지이고,
  $W_{m,n}(a,s)$와 $\gamma_{m,n}(a,s)$가 slip 장벽을 제공한다.
  외력을 포함한 최종식은 $\Phi(a,s,t)$이다.

  \item \textbf{히스테리시스 에너지 축적식:}
  관성 Markov 모델에서는 식~\eqref{eq:Kramers_dissipation}의 $\dot D_K$를 적분한다.
  식~\eqref{eq:overdamped_conditions}이 검증되고 경계가 no-flux 또는 periodic이면
  \[
    D_{\rm irr}(t)=\int_0^t\iint
    \left[
    \frac{J_a^2}{M_aP}+\frac{J_s^2}{M_sP}
    \right]\dd a\dd s\dd t'\ge0.
  \]
  흡수형 균열경계에서는 식~\eqref{eq:dissipation_with_boundary}의 경계유출항을 별도로
  포함한다.

  \item \textbf{평균거리식:}
  \[
    \bar a(t)=\iint aP(a,s,t)\dd a\dd s.
  \]

  \item \textbf{평균에너지식:}
  \[
    \bar E(t)=\iint\Phi(a,s,t)P(a,s,t)\dd a\dd s.
  \]

  \item \textbf{피로 및 균열개시를 나타내는 분포꼬리:}
  \[
    p_{\rm tail}(t;\sigma)=
    \int_{a^\ddagger(\sigma)}^\infty\int
    P(a,s,t)\dd s\dd a,
  \]
  또는 흡수경계를 통과한 누적 확률유속을 사용한다.
\end{enumerate}

\small
결론적으로 $K_\mu$가 들어간 급수는 임의의 hysteresis kernel이 아니라 무한 원자열의
정확한 에너지합을 빠르게 계산하기 위한 수학적 재표현이다. 실제 loading/unloading
비대칭은 이 에너지면 위에서 $\rho$ 또는 조건부로 축약된 $P$가 유한한 시간에
완화되면서 만드는 확률유속과,
필요한 경우 격자우물 사이의 비가역 $z$ 전이에서 발생한다. 순수 법선 모델은
피로개구와 분포꼬리 진화를 담당하고, $s,z$를 켠 확장모델은 이상적 단일슬립
소성을 담당한다. 알루미늄의 정량적 소성모델로 주장하기 전에는 FCC EAM/DFT
GSF와 전위경화변수를 추가 검증해야 한다.
\normalsize
\endgroup

\end{koreantranslation}

\end{document}
